% paper_18_exchange_constants.tex — GeoVac Paper 18
% Spectral-Geometric Exchange Constants: Projecting Discrete
% Spectra onto Continuous Manifolds
% Status: Draft, April 17, 2026
% Category: Core (papers/core/)
\documentclass[aps,pra,twocolumn,superscriptaddress,longbibliography]{revtex4-2}

\usepackage{amsmath,amssymb,amsthm,graphicx,xcolor,bm,braket,dcolumn,booktabs}
\usepackage{hyperref}

\newtheorem{theorem}{Theorem}
\newtheorem{observation}{Observation}
\newtheorem{conjecture}{Conjecture}

\begin{document}

\title{Spectral-Geometric Exchange Constants:\\
Projecting Discrete Spectra onto Continuous Manifolds}
\thanks{Paper 18 in the Geometric Vacuum series}

\author{J.~Loutey}
\affiliation{Independent Researcher, Kent, Washington}

\date{April 16, 2026}

%% ========================================================================
%% ABSTRACT
%% ========================================================================

\begin{abstract}
Compact Lie groups have discrete unitary representations
(Peter--Weyl theorem); the GeoVac graph exploits this discreteness
directly.  We show that the transcendental conversion factors
appearing throughout the GeoVac hierarchy---$\kappa = -1/16$
($S^3 \to \mathbb{R}^3$), the Stieltjes seed $e^a E_1(a)$
(Laguerre basis $\to$ prolate spheroid), the adiabatic eigenvalue
$\mu(R)$ ($S^3 \times S^3 \to$ hyperspherical)---are manifestations
of a single structural pattern: they arise when discrete
representations of compact groups are projected onto non-compact
coordinate systems used for measurement.  The taxonomy classifies
these factors by the type of compactness mismatch they mediate,
and connects them to the Weyl--Selberg exchange constants of
spectral geometry.  This identification explains why the constants
are transcendental (compact$\to$non-compact projection introduces
$\pi^{d/2}$), why their complexity increases with the geometry,
and why they vanish when the computation remains on the graph.
The falsifiable claim: the transcendental content of any GeoVac
observable is determined entirely by the type of projection
linking the compact discrete description to the non-compact
coordinate system.  Beyond the basis-level taxonomy, a
three-axis classification of QED transcendental content on~$S^3$
emerges: (1)~operator order (1st~vs 2nd) determines motivic weight
(odd~vs even~zeta); (2)~bundle type (scalar~vs spinor) modulates
coefficients; (3)~diagram topology introduces Dirichlet $L$-function
characters.  This classification is verified to be $R$-independent,
persisting in the flat-space limit via Weyl's law.
The taxonomy is identified as a classification of Mellin transforms
of heat kernels at successively higher perturbative orders: the free
heat kernel produces $\pi^{\mathrm{even}}$ (calibration tier), the
interaction-deformed kernel produces the full MZV algebra (embedding
tier), and operator order determines whether the Mellin integrand
samples squared or raw eigenvalues.
We note that the combination rule $K = \pi(B + F - \Delta)$ in the
conjectural fine structure constant formula (Paper~2) has the same
structural character: $\pi$ multiplying a sum of discrete
representation-theoretic data and continuous spectral invariants.
Phase-4G analysis further establishes that $B$, $F$, $\Delta$ have
categorically different origins (finite Casimir trace at $m=3$;
infinite Dirichlet series at $d_{\max}=4$; single-level Dirac-mode
count at $n=3$), so $K$ is a cross-sector sum of structurally
independent objects with no common generator. Finally, the RH-directed
sprint series (2026) produces three further concrete realizations of
the exchange-constant taxonomy: (a)~the closed-form identity
$D_{\mathrm{even}}(s) - D_{\mathrm{odd}}(s) = 2^{s-1}(\beta(s) -
\beta(s-2))$ for the Dirac Dirichlet difference, providing an exact
$\chi_{-4}$-content projection; (b)~the first GUE-like zero spacing
signature in the framework, on the spectral Dirichlet side only
(Ihara side is Berry-Tabor / Poisson), creating an "RMT-spacing
without critical-line confinement" tier; and (c)~a new
\emph{algebraic-but-non-radical} taxonomic tier realized by the
$S_6$-Galois factor $P_{12}$ of the scalar $S^5$ Ihara zeta at
$N_{\max}=3$ (12 zeros with no radical expression over~$\mathbb{Q}$).
Taken together, these instances exhibit a meta-pattern of
\emph{structural incommensurability within single framework
observables}: what appears to be one object in GeoVac repeatedly
resolves into several with categorically different origins. This
meta-pattern is consistent across three independently-derived
three-axis classifications and is discussed in Sec.~\ref{sec:meta_observation}.
\end{abstract}

\maketitle

%% ========================================================================
%% I. INTRODUCTION
%% ========================================================================

\section{Introduction}
\label{sec:intro}

The Geometric Vacuum program constructs discrete lattice
Hamiltonians by identifying the natural geometry of each quantum
system and discretizing it as a sparse
graph~\cite{loutey_paper0,loutey_paper1,loutey_paper7}.  A
persistent observation across the hierarchy of natural geometries
is that the discrete algebraic structures---integer eigenvalues,
finite recurrence relations, Gaunt integrals built from Wigner
$3j$ symbols---are entirely self-contained, but that their
connection to physical observables in continuous coordinate
systems always introduces transcendental constants.

The thesis of this paper is that these transcendental conversion
factors are consequences of \emph{compactness}.  Compact Lie groups
have discrete unitary representations (the Peter--Weyl
theorem~\cite{peter_weyl1927}); the GeoVac graph exploits this
discreteness directly.  Transcendental numbers enter when---and only
when---those discrete representations are translated into non-compact
coordinate systems used for measurement.  The taxonomy developed below
classifies the resulting conversion factors by the type of
compact-to-non-compact projection that produces them.  This yields a
falsifiable structural claim: the transcendental content of any GeoVac
observable is determined entirely by the type of projection linking
the compact discrete description to the non-compact coordinate system
(Claim~4, Sec.~\ref{sec:observable_classification}).

The simplest instance is already visible in Paper~0's packing
construction~\cite{loutey_paper0}.  Axiom~2 (geometric indifference)
forces states onto isotropic shells---circles, i.e.\ copies of~$S^1$.
The shell is compact; its circumference $2\pi r$ introduces $\pi$
into the area formula $\sigma_0 = \pi d_0^2/2$ and hence into the
state count $N_k = A_k / \sigma_0 = 2(2k-1)$.  The integer state count
is a consequence of compact topology; $\pi$ is the price of expressing
that count as a ratio of continuous areas.  This pattern---compact
geometry producing integer quantum numbers, with $\pi$ entering
through projection onto non-compact coordinates---recurs at every level
of the hierarchy.

The conversion factors are instances of Weyl's law~\cite{weyl1911} and
the Selberg trace formula~\cite{selberg1956}; what is new is their
systematic identification across the GeoVac hierarchy, culminating in a
falsifiable observable classification (Claim~4) and the observation that
operator order on the compact space determines the motivic weight of the
resulting transcendental content.


%% ========================================================================
%% II. CATALOG OF EXCHANGE CONSTANTS
%% ========================================================================

\section{Catalog of Exchange Constants}
\label{sec:catalog}

\subsection{Level 1: $\kappa = -1/16$ ($S^3 \to \mathbb{R}^3$)}
\label{sec:kappa}

The compactness of $S^3$ is what makes the graph eigenvalues discrete
integers; $\kappa$ converts this discrete spectrum to the non-compact
$\mathbb{R}^3$ energy scale.

The graph Laplacian on $S^3$ produces pure integer eigenvalues
\begin{equation}
  \lambda_n = -(n^2 - 1), \qquad n = 1, 2, 3, \ldots
  \label{eq:s3_eigenvalues}
\end{equation}
with degeneracy $g_n = n^2$~\cite{loutey_paper7}.  The physical
hydrogen energy levels $E_n = -1/(2n^2)$ are related by the
energy-shell constraint $p_0^2 = -2E$, which acts as a
stereographic focal length projecting $S^3$ onto flat
$\mathbb{R}^3$.  The universal kinetic scale factor
\begin{equation}
  \kappa = -\frac{1}{16}
  \label{eq:kappa}
\end{equation}
calibrates the ground state energy:
$\kappa \cdot \lambda_{\max} = E_H = -1/2$~Ha, where
$\lambda_{\max} \to 8$ is the maximum graph Laplacian eigenvalue
in the continuum limit~\cite{loutey_paper0}.  Per-shell energies
$E_n = -Z^2/(2n^2)$ emerge from the full graph Hamiltonian
$H = \kappa \mathcal{L} + W$ (with node weights $W_{nn} = -Z/n^2$),
not from $\kappa$ alone.

The origin of $\kappa$ is well understood geometrically: it is
the Jacobian of the stereographic projection composed with the
conformal factor of the Fock map.  Specifically, if $\Omega$
is the conformal factor $\Omega = (1 + p^2/p_0^2)^{-1}$, then
the Laplace--Beltrami operator on $S^3$ and the flat-space
kinetic operator are related by
$\nabla^2_{S^3} = \Omega^{-2}[\nabla^2_{\mathrm{flat}} + \text{lower order}]$,
and $\kappa$ absorbs the accumulated conformal rescaling and
the $p_0$ energy-shell normalization~\cite{loutey_paper7}.

The key property of $\kappa$ is that it is \textit{unnecessary
on the graph}.  The graph Laplacian eigenvalue problem $L\psi =
\lambda\psi$ is solved entirely in terms of integers.  The
constant $\kappa$ appears only when the result is translated
into atomic units (Hartrees, bohrs)---i.e., when the discrete
spectrum is projected onto a continuous coordinate system.

\subsection{Level 2: $e^a E_1(a)$ (Laguerre basis $\to$ prolate spheroid)}
\label{sec:stieltjes}

The Laguerre spectral basis inherits the compact angular structure of
$S^3$; the Stieltjes seed $e^a E_1(a)$ enters when this basis is
embedded in the non-compact prolate spheroidal coordinate, whose
curvature introduces a pole absent from the compact description.

The prolate spheroidal radial solver uses a Laguerre spectral
basis $\varphi_n(\xi) = e^{-\alpha(\xi-1)} L_n(2\alpha(\xi-1))$
with matrix elements computed from three-term
recurrences~\cite{loutey_paper11}.  For $\sigma$ states ($m = 0$),
all matrix elements are purely algebraic: the Laguerre moment
matrices $M_k[i,j] = \int_0^\infty x^k L_i(x) L_j(x) e^{-x}
dx$ are band-structured matrices with rational entries, built
from the recurrence $x L_n = -(n+1)L_{n+1} + (2n+1)L_n -
nL_{n-1}$.

For $m \neq 0$ states, the centrifugal term $m^2/(\xi^2 - 1)$
introduces a $1/x$ singularity in the Laguerre variable
$x = 2\alpha(\xi - 1)$.  The associated Laguerre basis
$L_n^{|m|}(x)$ with weight $x^{|m|} e^{-x}$ absorbs this
singularity by construction.  The centrifugal matrix elements
decompose via partial-fraction expansion:
\begin{equation}
  \frac{m^2}{\xi^2 - 1}
  = 4\alpha^2 m^2 \left[\frac{1}{x}
    - \frac{1}{x + 4\beta}\right],
  \label{eq:partial_fraction}
\end{equation}
where $\beta = \alpha(\xi_0 - 1)$ is the Laguerre scale offset.
The two terms have fundamentally different algebraic character:
\begin{itemize}
  \item The $1/x$ term produces a lowered moment
    $M_{-1}[i,j] = \int_0^\infty x^{|m|-1} L_i^{|m|}(x)
    L_j^{|m|}(x)\,e^{-x}\,dx$, which is fully algebraic via
    the DLMF 18.9.13 summation
    identity~\cite{dlmf_laguerre}.
  \item The $1/(x + 4\beta)$ term produces a Stieltjes matrix
    $J[i,j] = \int_0^\infty \frac{x^{|m|}}{x + a}\,
    L_i^{|m|}(x) L_j^{|m|}(x)\,e^{-x}\,dx$, where $a = 4\beta$.
    This matrix satisfies a three-term recurrence in the
    row/column indices, seeded by the single transcendental
\end{itemize}
\begin{equation}
  S_0^{(\alpha)}(a) = \Gamma(\alpha+1)\,e^a\,E_{\alpha+1}(a),
  \label{eq:stieltjes_seed}
\end{equation}
where $E_n(a) = \int_1^\infty t^{-n} e^{-at}\,dt$ is the
generalized exponential integral and $a = 4\beta$ is the
Laguerre scale parameter~\cite{loutey_track_j}.  For integer
$\alpha = |m|$, all higher $E_n$ reduce to $E_1(a)$ via a
finite rational recurrence, so the entire matrix is built from
algebraic operations plus one transcendental constant.

The computational cost of the transcendental seed is negligible:
one call to $E_1(a)$, which achieves machine precision for
$a \ge 4$ (covering all physical parameters).  The seed then
propagates through an $O(N^2)$ algebraic recurrence that produces
all matrix elements.  The continuous mathematics has been
compressed from a \textit{function} (the quadrature integrand
over $[0,\infty)$) to a \textit{constant} (the Stieltjes seed),
with the entire function-level information reconstructed
algebraically from that constant.

An unexpected computational bonus: the associated Laguerre basis
converges \textit{faster} than the ordinary Laguerre basis for
$m \neq 0$ states.  For $|m| = 1$ ($\pi$ states), the associated
basis is stable by $N = 10$, compared to non-monotonic convergence
at $N \approx 20$ for the ordinary basis.  For $|m| = 2$
($\delta$ states), algebraic-quadrature agreement reaches
$\sim 5 \times 10^{-9}$ (machine precision), with PES shape
(equilibrium geometry) preserved exactly~\cite{loutey_track_j}.
This suggests that identifying the exchange constant and absorbing
it into the basis does not merely explain the transcendence---it
\textit{improves the spectral method}, because the adapted basis
eliminates the curvature mismatch that the exchange constant
mediates.

The Stieltjes seed $e^a E_1(a)$ is the Laplace transform of the
rational function $1/(1+t)$ evaluated at $a$.  It encodes how
the Laguerre inner product (defined on the flat half-line
$[0, \infty)$ with exponential weight) ``sees'' a shifted pole
at $x = -a$ arising from the curvature of the prolate spheroidal
coordinate system.  On the graph, there is no pole; the
centrifugal barrier is a selection rule encoded in the quantum
numbers.  The transcendental appears only when this selection
rule is represented as a rational function in a continuous
coordinate.

\subsection{Level 3: $\mu(R)$ ($S^3 \times S^3 \to$ hyperspherical)}
\label{sec:mu}

On the compact product space $S^3 \times S^3$, the two-electron
Hamiltonian has purely algebraic matrix elements; the adiabatic
eigenvalue $\mu(R)$ enters only when this compact description is
reparameterized into the non-compact hyperspherical coordinate~$R$.

The two-electron problem on $S^3 \times S^3$ is fully algebraic:
the Hamiltonian matrix has entries given by Gaunt integrals
(finite sums of Wigner $3j$ symbols with rational coefficients),
and the eigenvalues are roots of a polynomial with algebraic
coefficients.  The Full CI solver exploits this, achieving
0.35\% error for helium with no transcendentals
anywhere~\cite{loutey_fci}.

The hyperspherical reparameterization introduces coordinates
$(R, \alpha)$ where $R = \sqrt{r_1^2 + r_2^2}$ is the
hyperradius and $\alpha = \arctan(r_2/r_1)$ is the hyperangle.
The angular eigenvalues $\mu(R)$ depend parametrically on $R$
because the effective angular Hamiltonian is
$H(R) = H_0 + R \cdot V^{\mathrm{coupling}}$~\cite{loutey_paper13}.

Each perturbation coefficient $a_k$ in the expansion
$\mu(R) = \mu^{(0)} + a_1 R + a_2 R^2 + \cdots$ has a closed
form: the $a_k$ are finite sums of algebraic numbers (square
roots of rationals from Clebsch--Gordan coefficients) times
powers of $\pi^{-1}$ (from normalization over the angular domain
$[0, \pi/2]$), arising from partial harmonic sums over
$\mathrm{SO}(6)$ Casimir eigenstates~\cite{loutey_paper13}.
For example, $a_1(n{=}1, Z) = (8/3\pi)(\sqrt{2} - 4Z)$.
However, the full function $\mu(R)$ is not transcendental in the
sense originally claimed.  The angular Hamiltonian
$H(R) = H_0 + R \cdot V_C$ is a linear matrix pencil, so its
eigenvalues satisfy the characteristic polynomial
\begin{equation}
  P(R, \mu) \;=\; \det\!\bigl(H_0 + R\,V_C - \mu\,I\bigr) = 0,
  \label{eq:char_poly}
\end{equation}
which is a polynomial of degree $d$ in both $R$ and $\mu$ (where
$d$ is the matrix dimension at a given truncation).  The coefficients
of $P$ lie in $\mathbb{Q}(\pi, \sqrt{2})$, since $H_0$ has integer
entries (SO(6) Casimir eigenvalues) and $V_C$ has entries involving
$\pi^{-1}$ and $\sqrt{2}/\pi$ from the angular integrals.  This
makes $\mu(R)$ an \emph{algebraic} function of $R$ over the
coefficient ring $\mathbb{Q}(\pi, \sqrt{2})$---it is a root of
a polynomial, not a transcendental
function~\cite{loutey_track_p1}.

The transcendental content of $\mu(R)$ is therefore
\emph{inherited} from the coupling matrix $V_C$, which contains
upstream exchange constants of embedding type ($\pi^{-1}$ from
angular normalization, $\sqrt{2}$ from Gaunt-type integrals).
The $R$-dependence itself introduces no new transcendentals:
it is purely algebraic.

The failure of the Taylor/Pad\'e expansion at large $R$
(Track~G~\cite{loutey_track_g}) is reinterpreted as a
\emph{branch point} phenomenon, not evidence of transcendence.
The algebraic curve $P(R, \mu) = 0$ has branch points where
eigenvalue sheets meet---these are the zeros of the discriminant
$\Delta(R) = \mathrm{disc}_\mu(P)$.  The Taylor series around
$R = 0$ has a convergence radius set by the nearest branch point
in the complex $R$-plane.  Pad\'e approximants, being rational
functions, also cannot cross branch points where $\mu(R)$ has
square-root type singularities.  But the implicit relation
$P(R, \mu) = 0$ is exact for \emph{all} $R$, including $R > 5$
where the series diverges.  This has been verified computationally
to machine precision at $R = 50$~bohr~\cite{loutey_track_p1}.

The $O(R) \to O(R^2)$ crossover, which was previously interpreted
as a regime change requiring transcendental interpolation, is
simply the behavior of an algebraic function whose dominant term
changes character: at small $R$, the linear term $V_C \cdot R$
dominates; at large $R$, higher-order terms in the characteristic
polynomial dominate.  The crossover is smooth along the algebraic
curve, even though it is singular from the perspective of a Taylor
expansion centered at $R = 0$.

The crucial point remains that $\mu(R)$ is not intrinsic to the
two-electron problem.  It is an artifact of the hyperspherical
coordinate system.  The $S^3 \times S^3$ product space does not
require $\mu(R)$; it only appears when one reparameterizes for
computational efficiency.  The exchange rate between the discrete
product-space description and the continuous hyperspherical
description is an algebraic function of $R$, with transcendental
content inherited from the angular coupling integrals.


\subsection{Level 4: $\mu(\rho,R)$ (product space $\to$ molecule-frame hyperspherical)}
\label{sec:mu_level4}

At Level~4 the compact product basis (prolate spheroidal CI) is fully
algebraic but structurally incomplete; the non-compact mol-frame
hyperspherical reparameterization introduces $\mu(\rho,R)$ as the
irreducible cost of accessing three-body coalescence physics that the
compact description cannot represent.

Level~4 combines the two-center complication of Level~2 with the
two-electron complication of Level~3.  Two approaches exist, with
fundamentally different algebraic character.

The product-space approach (prolate spheroidal CI, Paper~12) is
fully algebraic: all electron-electron matrix elements are computed
from Neumann expansion recurrence relations with zero spatial
quadrature.  The Laguerre radial and Legendre angular matrix
elements are likewise algebraic.  No transcendental constant
appears anywhere in the computation.

However, a systematic convergence study (Track~M) demonstrates
that this algebraic approach saturates at $\sim$92.5\% of the
exact dissociation energy $D_e$~\cite{loutey_track_m}.  The
angular basis convergence is rapid---92.3\% at $l_{\max} = 2$
(46 basis functions), 92.5\% at $l_{\max} = 4$ (114 basis
functions)---but successive improvements decay by a factor of
200 from $l = 1 \to 2$ to $l = 2 \to 3$, with the $l = 3 \to 4$
increment contributing only 0.04\%.  Extrapolation in the slow
regime predicts that 5\% error would require $l_{\max} \sim 45$
and 1\% error $l_{\max} \sim 113$, both impractical.  The
effective ceiling is 92.5\% $D_e$.

The smoking gun is explicit correlation: 9~basis functions with
$r_{12}^2$ factors (Hylleraas $p = 2$) achieve 94.7\% $D_e$,
exceeding the 114-function prolate CI plateau.  The residual
7.5\% matches Paper~12's cusp diagnosis (7.6\%): the
electron-electron cusp ($1/r_{12}$ singularity) is a completeness
deficiency of the product basis, not a convergence issue.
The transcorrelated approach (Paper~14, Sec.~IV; Paper~15,
Sec.~VI.J) removes this singularity via a Jastrow similarity
transformation, converting $1/r_{12}$ from a hard embedding
constant to a smooth gradient operator at the cost of
non-Hermiticity (Sec.~\ref{sec:taxonomy}).

The mol-frame hyperspherical reparameterization (Paper~15)
resolves this by introducing coordinates $(\rho, R, \alpha)$
where $\rho$ is the electron hyperradius and $\alpha$ the
hyperangle.  The angular eigenvalues $\mu(\rho, R)$ are now a
transcendental function of \textit{two} continuous parameters,
reflecting the simultaneous three-body coalescence in coordinate
and configuration space.  This solver achieves 94.1\% $D_e$,
surpassing the prolate CI ceiling.

The distinction between Levels~3 and 4 is qualitative.  At
Level~3, the $S^3 \times S^3$ FCI converges (albeit slowly)---the
exchange constant $\mu(R)$ is a computational convenience that
accelerates convergence but is not fundamentally necessary.  At
Level~4, the prolate product-space CI \textit{saturates}: the
exchange constant $\mu(\rho, R)$ is irreducible because the
product basis is structurally blind to three-body coalescence.
The reparameterization does not merely accelerate convergence;
it accesses physics that the algebraic product space cannot
represent at any practical basis size.


%% ========================================================================
%% III. THE PATTERN
%% ========================================================================

\subsection{Level 5: RH-sprint exchange constants (2026)}
\label{sec:level5_rh_sprint}

The 2026 RH-directed sprint series added four distinct concrete
realizations of the exchange-constant taxonomy to the catalog.
All four live on the same S^3 substrate as Levels 1--4 but classify
different projections of the spectral data.

\paragraph{Closed-form $\chi_{-4}$ content in the Dirac Dirichlet
difference (RH-J).}  On $S^3$ with the Camporesi--Higuchi Dirac
spectrum $|\lambda_n| = n + 3/2$ and degeneracy $g_n = 2(n+1)(n+2)$,
define the Dirac Dirichlet series $D(s) = \sum_n g_n / |\lambda_n|^s$
and its vertex-parity split $D_{\mathrm{even}}(s)$, $D_{\mathrm{odd}}(s)$
(Paper~28~\cite{loutey_paper28}, eqs.~(3)--(5)). The difference has
the exact closed form
\begin{equation}
  D_{\mathrm{even}}(s) - D_{\mathrm{odd}}(s) \;=\; 2^{s-1}\,\bigl(
    \beta(s) - \beta(s-2)\bigr), \quad s \in \mathbb{Z}_{\ge 2},
  \label{eq:rh_j_identity}
\end{equation}
where $\beta(s) = L(s, \chi_{-4})$ is the Dirichlet $L$-function of
the mod-4 character. Two elementary Hurwitz identities drive the
collapse:
$\zeta(s, 5/4) = \zeta(s, 1/4) - 4^s$ and
$\zeta(s, 3/4) = \zeta(s, 1/4) - 4^s \beta(s)$.
Verified by sympy symbolic identity at $s \in \{2, \ldots, 10, 15\}$
and by PSLQ at 100-digit precision
(\texttt{debug/spectral\_chi\_neg4\_memo.md}).
Eq.~(\ref{eq:rh_j_identity}) is the concrete realization of
Paper~28's qualitative "vertex parity acts as $\chi_{-4}$" statement:
the Dirichlet $L$-content enters the spectral-Dirichlet difference
as a closed-form linear combination of $\beta(s)$ and $\beta(s-2)$
with integer prefactor $2^{s-1}$.

\paragraph{RMT-spacing without critical-line confinement: the
spectral-Ihara dichotomy (RH-M, RH-G).}  The zeros of $D(s)$,
$D_{\mathrm{even}}(s)$, $D_{\mathrm{odd}}(s)$ in the strip
$\mathrm{Im}(s) \in [0, 70]$ have unfolded imaginary-part spacing
coefficient of variation $\mathrm{CV} \in \{0.348, 0.343, 0.396\}$
---consistent with GUE ($\mathrm{CV} \approx 0.46$) and inconsistent
with Poisson ($\mathrm{CV} \approx 1.05$). Kolmogorov--Smirnov testing
rejects Poisson at $p = 0.004, 0.030$ respectively and cannot reject
GUE. This is the first GUE signature in the GeoVac framework, and
it is confined to the \emph{spectral} Dirichlet side only: the
combinatorial / Ihara-zeta side (graph non-backtracking walks) has
Poisson / Berry-Tabor spacing at $\mathrm{CV} \approx 1$, a categorical
dichotomy. Equally important, the $D(s)$ zeros are \emph{not} on a
single critical line (they are scattered in
$\mathrm{Re}(s) \in [2.02, 3.42]$), so the GUE signature is a
spacing-only statement, not an RH-style line constraint: classical
$\zeta$ has \emph{both} RMT spacing and critical-line confinement;
GeoVac's spectral zeta has only the first. This creates a new
taxonomic tier, \emph{RMT-spacing without critical-line confinement}.

\paragraph{Algebraic-but-non-radical transcendental content (RH-F).}
The integer-coefficient Ihara-zeta factorizations produced in
Paper~29~\cite{loutey_paper29} factor over specific number fields as
follows. Cyclotomic factors $\Phi_3$, $\Phi_4$, $\Phi_6$ have Galois
group $\mathbb{Z}/2$ (abelian); they fingerprint specific graph cycles
(e.g.~$\Phi_3, \Phi_6 \leftrightarrow$ 3-prism cycles of the $\ell=1$
Coulomb $S^3$ component). Ramanujan-boundary factors
$(4s^2+1)^2$ and $(9s^2+1)^4$ live in $\mathbb{Q}(i)$, consistent
with the Hopf $\mathbb{Z}_2$ $m$-reflection sub-action of the Paper~25
Hopf $U(1)$ (see Observation~2 of~\cite{loutey_paper29}).
\emph{Bulk} factors have non-abelian Galois groups: $S_3$ for scalar
$S^3$ cubics; $S_4$ for the Dirac-$S^3$ Rule-A $n_{\max}=3$ quartic;
and crucially $S_6$ (non-solvable) for the factor $P_{12}(s^2)$ in
the scalar $S^5$ $N_{\max}=3$ Ihara zeta. By Abel--Ruffini, the 12
zeros of $P_{12}$ have \emph{no radical expression} over
$\mathbb{Q}$; by Kronecker--Weber, they do not lie in any cyclotomic
extension $\mathbb{Q}(\zeta_M)$. This places a new tier in
Sec.~\ref{sec:taxonomy}'s output-class axis: \emph{algebraic but
not radical} --- intermediate between closed-form radical algebraic
and classical periods.

\paragraph{$\alpha^2$-weighted Ihara zeta in the ring
$\mathbb{Q}(\alpha^2)[s]$ (RH-K).}  The unweighted Ihara zeta
depends only on $0/1$ graph connectivity. Introducing the edge
weighting $w(a, b) = 1 + \alpha^2 \cdot |f_{\mathrm{SO}}(a) +
f_{\mathrm{SO}}(b)|/2$, where
$f_{\mathrm{SO}}(n, \kappa) = -(\kappa+1)/[4 n^3 \ell(\ell+1/2)(\ell+1)]$
is the dimensionless Breit--Pauli spin-orbit diagonal, gives a
weighted Ihara zeta $\zeta_G^w(s)$ living in the ring
$\mathbb{Q}(\alpha^2)[s]$ --- an explicit realization of the
spinor-intrinsic tier without the $\gamma = \sqrt{1 - (Z\alpha)^2}$
radical (\texttt{geovac/ihara\_zeta\_weighted.py}, 14/14 tests).
For Rule-A at $n_{\max}=3$, the inverse zeta expands as a polynomial
in $\alpha$ of degree 44, with $\alpha^2$ coefficient carrying
the clean rational prefactor $-35/648$. The coefficient is a
graph-invariant, not an edge-local perturbation --- a concrete
$\alpha^2$-content projection.

\subsection*{Taxonomic placement of the Level-5 realizations}

The four Level-5 constants slot into the existing taxonomy as
follows. The $\chi_{-4}$ identity (\ref{eq:rh_j_identity}) is a
\emph{flow} exchange constant (infinite spectral sum), with the
novelty that its value is an integer-coefficient linear combination
of classical Dirichlet $L$-values, so Paper~18's transcendental
taxonomy gains a "Dirichlet-$L$-closed" cell at the intersection of
flow and Dirac-spectrum projection. The GUE-vs-Poisson dichotomy is
a \emph{meta-level} discriminant, not itself a transcendental, but
classifies spectral objects by zero statistics --- a new axis in
Sec.~\ref{sec:taxonomy}. The algebraic-but-non-radical tier from
$P_{12}$ is a new \emph{intrinsic} exchange-constant sub-class
between Level~0 (rational/integer) and classical periods; it is the
first time a non-solvable Galois group has been exhibited as a
framework-native transcendental class. The $\alpha^2$-weighted
zeta is a \emph{spinor-intrinsic} exchange constant in the ring
$\mathbb{Q}(\alpha^2)[s]$, a subring of the full spinor-intrinsic
tier $\mathbb{Q}(\alpha^2)[\gamma]/(\gamma^2 + (Z\alpha)^2 - 1)$
introduced in Sec.~\ref{sec:taxonomy}.

% ------------------------------------------------------------------

\section{The Pattern: Weyl--Selberg Exchange}
\label{sec:pattern}

\subsection{Weyl's law}

The mathematical content of the observations in
Sec.~\ref{sec:catalog} is an instance of Weyl's law, which
states that the eigenvalue counting function $N(\lambda)$ of
the Laplace--Beltrami operator on a compact Riemannian manifold
of dimension $d$ satisfies~\cite{weyl1911}
\begin{equation}
  N(\lambda) \sim \frac{\omega_d}{(2\pi)^d}\,
  \mathrm{vol}(\Omega)\,\lambda^{d/2}
  \qquad (\lambda \to \infty),
  \label{eq:weyl}
\end{equation}
where $\omega_d = \pi^{d/2}/\Gamma(d/2+1)$ is the volume of
the unit ball in $\mathbb{R}^d$.

Weyl's law relates two fundamentally different descriptions
of the same manifold:
\begin{itemize}
  \item The \textbf{spectral side}: $N(\lambda)$, a discrete
    staircase function counting integer-labeled eigenvalues.
  \item The \textbf{geometric side}: $\mathrm{vol}(\Omega)$, a
    continuous quantity measured by integration over the manifold.
\end{itemize}
The conversion factor $\omega_d / (2\pi)^d$ involves
$\pi^{d/2}$---a transcendental that depends only on the
dimension.  This factor is the prototypical
\textit{spectral-geometric exchange constant}: it converts
between counting (discrete) and measuring (continuous).

For the specific manifolds in the GeoVac hierarchy, the ball
volumes $\omega_d$ carry $\pi^{\lfloor d/2 \rfloor}$:
\begin{itemize}
  \item $S^1$ ($d = 1$): $\omega_1 = 2$,
    Weyl constant $\omega_1/(2\pi)^1 = 1/\pi$.
  \item $S^2$ ($d = 2$): $\omega_2 = \pi$,
    Weyl constant $\omega_2/(2\pi)^2 = 1/(4\pi)$.
  \item $S^3$ ($d = 3$): $\omega_3 = 4\pi/3$,
    Weyl constant $\omega_3/(2\pi)^3 = 1/(6\pi^2)$.
  \item $S^5$ ($d = 5$): $\omega_5 = 8\pi^2/15$,
    Weyl constant $\omega_5/(2\pi)^5 = 1/(60\pi^3)$.
\end{itemize}
In each case, the Weyl constant involves inverse powers of $\pi$
(specifically $\pi^{-\lceil d/2 \rceil}$), reflecting the
transcendental cost of converting a discrete eigenvalue count
into a continuous volume measure.

\subsection{Compactness as the source of discreteness}

The Peter--Weyl theorem~\cite{peter_weyl1927} states that the unitary
representations of a compact Lie group form a discrete, countable,
orthogonal decomposition of $L^2(G)$.  This is why quantum numbers are
integers: $\mathrm{SO}(4)$ is compact, so its irreps are labeled by
discrete indices $(n, l, m)$ with finite multiplicities $g_n = n^2$.
Exchange constants arise when these discrete representations are
expressed in non-compact coordinate systems---the stereographic
$S^3 \to \mathbb{R}^3$ introduces $\kappa$; embedding in prolate
spheroidal coordinates introduces $e^a E_1(a)$; reparameterizing
$S^3 \times S^3$ as hyperspherical introduces $\mu(R)$.  The
transcendental content is the toll for crossing the compactness
boundary.

\subsection{$\pi$ as the founding example}

The simplest spectral-geometric exchange constant is $\pi$ itself.
The Laplacian on $S^1$ (the circle) has eigenvalues $n^2$, counted
by integers; the continuous circumference is $2\pi$.  The
conversion between the discrete eigenvalue count and the continuous
measure is mediated by $\pi$.  More generally, the volume of the
unit $d$-sphere involves $\pi^{d/2}$, and the Weyl constant
$\omega_d / (2\pi)^d$ carries exactly this factor.  Every exchange
constant in the GeoVac hierarchy is a descendant of this primordial
example: $\kappa$ inherits from the $S^3$ Weyl constant, the
Stieltjes seed involves the exponential integral (a function whose
Taylor coefficients involve the Euler--Mascheroni constant, itself
a limit of harmonic sums weighted by $\pi$-dependent terms), and
$\mu(R)$ is an eigenvalue of a Hamiltonian on $S^5$ whose Weyl
asymptotics carry $\pi^{5/2}$.

\begin{observation}[The graph is $\pi$-free]
\label{obs:pi-free}
Every algebraic operation in the GeoVac graph construction
produces integers or rationals:
\begin{itemize}
  \item Eigenvalues of the graph Laplacian: $\lambda_n = -(n^2 - 1)$,
    integer for all $n$.
  \item Degeneracies: $g_n = n^2$, integer.
  \item Gaunt integrals (Legendre normalization,
    $\int P_{l_1} P_k P_{l_2}\,dx$): ratios of factorials via
    Wigner $3j$ symbols, rational.
  \item Selection rules: combinatorial (triangle inequality on angular
    momentum labels).
  \item Clebsch--Gordan coefficients: algebraic numbers (square roots
    of rationals).
  \item Casimir eigenvalues: $l(l+1)$, $\nu(\nu+4)/2$, integer or
    half-integer for all levels.
\end{itemize}
Transcendental numbers ($\pi$, $e^a E_1(a)$, $\zeta(2)$) enter
exclusively through normalization to continuous measures:
$\int |Y|^2\,d\Omega = 1$ (with $d\Omega$ carrying $4\pi$),
$\mathrm{vol}(S^3) = 2\pi^2$, and spectral zeta functions.
The exchange constants are precisely the $\pi$-injection points.
\end{observation}

\subsection{The Selberg trace formula}

Weyl's law is asymptotic.  The Selberg trace
formula~\cite{selberg1956} provides the exact identity:
\begin{equation}
  \sum_j h(r_j)
  = \frac{\mathrm{Area}(\Gamma \backslash \mathcal{H})}{4\pi}
    \int_{-\infty}^{\infty} r\,\tanh(\pi r)\,h(r)\,dr
    + \sum_{\{\gamma\}} \cdots
  \label{eq:selberg}
\end{equation}
where the left side sums over the discrete spectrum, the first
term on the right involves the continuous volume (with a factor
of $\pi$), and subsequent terms sum over conjugacy classes
(geodesics).  The Selberg trace formula is the bridge between
the spectral world and the geometric world, and the
transcendental factors ($\pi$, $\tanh$, logarithms of
algebraic numbers) are the toll for crossing.

When the quotient $\Gamma \backslash G$ is compact, the spectrum
is purely discrete and the trace formula involves only algebraic
data plus $\pi$-dependent conversion factors.  When the quotient
is non-compact, a continuous spectrum appears and the exchange
involves Eisenstein series~\cite{selberg1956}.  This parallels the
GeoVac hierarchy: Level~1 (compact $S^3$) requires only a constant
($\kappa$); Level~3 involves an algebraic function $\mu(R)$ with
transcendental coefficients (Sec.~\ref{sec:algebraic_curve}).

\subsection{The Mellin transform as the taxonomic engine}
\label{sec:mellin}

The preceding subsections identify the \emph{what}---compact-to-non-compact
projection introduces transcendental content---and the \emph{where}---Weyl's
law and the Selberg trace formula are the structural bridges.  Here we
identify the \emph{how}: the Mellin transform of the heat kernel trace is
the single mathematical operation that converts the discrete spectral data
of a compact manifold into the transcendental content classified in
Sec.~\ref{sec:taxonomy}.

For a non-negative self-adjoint operator $A$ on a compact Riemannian
manifold, the spectral zeta function is related to the heat kernel trace
by~\cite{seeley1967,minakshisundaram1949}
\begin{equation}\label{eq:mellin_bridge}
  \zeta_A(s) \;=\; \frac{1}{\Gamma(s)} \int_0^\infty t^{s-1}\,
  \operatorname{Tr}\bigl(e^{-tA}\bigr)\,dt\,,
\end{equation}
where $\operatorname{Tr}(e^{-tA}) = \sum_n g_n\,e^{-t\lambda_n}$ is a
purely spectral object: the sum runs over the discrete eigenvalues
$\lambda_n$ with their integer degeneracies~$g_n$.  Both sides of
Eq.~\eqref{eq:mellin_bridge} encode the same information---the spectrum
of~$A$---but on the left in a form that produces transcendental numbers
($\pi^{d/2}$ from the Seeley--DeWitt asymptotic expansion of the heat
trace, odd-zeta values from the analytic continuation), and on the right
in a form that is manifestly a function of integer-labeled eigenvalues
and degeneracies.  The Mellin transform is the bridge; the $\Gamma(s)$
prefactor is the conversion toll.

\paragraph{Operator order determines the Mellin output class.}
The three-axis classification of QED transcendentals on~$S^3$
(Sec.~\ref{sec:taxonomy}) has its foundation in the distinction between
applying the Mellin transform to the squared eigenvalues $\lambda_n^2$
versus the raw eigenvalues $|\lambda_n|$.

For the \emph{squared} Dirac operator $D^2$ on unit~$S^3$, the
eigenvalues are $(n + 3/2)^2$, which are perfect squares of half-integers.
The heat trace $\operatorname{Tr}(e^{-tD^2})$ is a theta-function-like
object whose Mellin transform samples only even powers: the variable
substitution $m = 2n+3$ converts the spectral sum to a sum over odd
integers, and the Euler--Bernoulli formula for $\zeta_R(2k)$ maps every
term to $\text{rational} \times \pi^{2k}$.  This is the content of the
T9 theorem (Paper~28, Theorem~1~\cite{loutey_paper28}):
$\zeta_{D^2}(s)$ is a two-term polynomial in $\pi^2$ at every
integer~$s$.  No odd-zeta value can appear because the squaring
operation has removed the odd-power structure from the Mellin integrand.

For the \emph{first-order} Dirac operator $|D|$ with eigenvalues
$n + 3/2$, the Mellin transform accesses the raw half-integer spectrum.
The Dirac Dirichlet series $D(s) = \sum g_n |\lambda_n|^{-s}$ involves
both even and odd powers of the spectral parameter, and the parity
discriminant (Paper~28, Theorem~2~\cite{loutey_paper28}) follows
directly: at even~$s$, both constituent Hurwitz zeta terms are
Bernoulli-type ($\pi^{\text{even}}$); at odd~$s$, both involve
genuinely independent odd-zeta values.  The operator order---first
versus second---determines whether the Mellin transform samples the
raw spectrum or its square, and this single distinction generates the
first axis of the three-axis classification.

\paragraph{Interaction deforms the heat kernel.}
The Coulomb potential $V = -Z/r$ perturbs the free Laplacian on~$S^3$.
In the heat kernel framework, the perturbed trace
$\operatorname{Tr}(e^{-t(A + V)})$ has an asymptotic expansion whose
Seeley--DeWitt coefficients $a_k(A + V)$ differ from the free
coefficients $a_k(A)$ by terms involving the potential and its
derivatives~\cite{gilkey1975}.  The Mellin transform of this perturbed
heat trace produces spectral zeta values that access transcendental
content beyond the $\pi$-polynomial ring of the free operator.
(For the \emph{free} Dirac operator~$D$ on unit~$S^3$, the SD
expansion terminates at exactly two terms by the Jacobi theta modular
identity: $\operatorname{Tr} e^{-tD^2}
= (\sqrt{\pi}/2)\,t^{-3/2} - (\sqrt{\pi}/4)\,t^{-1/2}
+ O(e^{-\pi^2/t})$, with $a_k = 0$ for $k \geq 2$; interaction opens
the full MZV algebra.)

The Sommerfeld fine-structure sums $D_p = \sum_{n=1}^\infty c_p(n)$
(Paper~28, Sec.~7~\cite{loutey_paper28}) are precisely the
perturbative expansion of this effect.  At each order $p$ in
$(Z\alpha)^{2p}$, the coefficients $c_p(n)$ involve generalized
harmonic numbers $H_{n-1}^{(r)}$, which are partial Euler sums that
access the full Euler--Zagier multiple zeta value (MZV) algebra.  The
free heat kernel produces only $\pi^{\text{even}}$ (Bernoulli); the
Coulomb interaction opens the door to odd-zeta values ($\zeta_R(3)$
at $p = 2$, $\zeta_R(5)$ and $\zeta_R(7)$ at $p = 3$, and so on)
and to products of zeta values at higher orders.

This maps directly onto the exchange constant taxonomy:
\begin{itemize}
  \item \textbf{Intrinsic tier} (angular/Gaunt structure): no Mellin
    transform is needed; the result is purely rational.
  \item \textbf{Calibration tier} ($\kappa$, $\pi$, $\zeta_R(2)$):
    Mellin transform of the \emph{free} heat kernel on the compact
    manifold $\to$ $\pi^{\text{even}}$ exclusively.
  \item \textbf{Embedding tier} (cusp, $1/r_{12}$, Sommerfeld $D_p$):
    Mellin transform of the \emph{interaction-deformed} heat kernel
    $\to$ the full MZV algebra.
\end{itemize}

\paragraph{The product survival rule as a Mellin selection rule.}
The product survival rule (Paper~28~\cite{loutey_paper28})---the systematic cancellation
of $\zeta(2) \times \zeta(\text{odd})$ products at every Sommerfeld
order---has a natural interpretation in this framework.  The
$n^2$ Fock degeneracy that weights the Sommerfeld sums is the same
$n^2$ that produces $F = \zeta_R(2) = D_{n^2}(4)$ in the free-geometry
Mellin output (Paper~2~\cite{loutey_paper2}, Phase~4F identification).
When the interaction-deformed Mellin transform generates product
terms $\zeta(a) \times \zeta(b)$, the Fock-degeneracy weighting
acts as a selection rule: products involving $\zeta(2)$ are
systematically expelled because $\zeta(2)$ is already ``used'' as
the Fock Dirichlet evaluation $F$, and the $n^2$ weighting enforces
orthogonality between the free and interacting Mellin outputs.

\paragraph{K--Sommerfeld separation as a Mellin consequence.}
The K--Sommerfeld structural separation
(Paper~28~\cite{loutey_paper28})---the PSLQ-verified fact that
$K/\pi$ is not in the $\mathbb{Q}$-span of $\{D_2, \ldots, D_6, 1\}$%
---is a direct consequence of the Mellin framework.
The ingredients of $K = \pi(B + F - \Delta)$ are all
\emph{free-geometry} Mellin invariants: $B$ is a truncated
trace of the free Casimir, $F = \zeta_R(2)$ is the spectral zeta
of the free Fock-degeneracy Dirichlet series at $s = d_{\max}$,
and $\Delta^{-1} = g_3^{\text{Dirac}}$ is a single-level mode
count of the free Dirac operator.  All three live in the
$\pi$-polynomial ring guaranteed by Eq.~\eqref{eq:zeta_d2}.
The Sommerfeld $D_p$, by contrast, are Mellin transforms of the
\emph{Coulomb-perturbed} heat kernel, dressed with harmonic
numbers that access the full MZV algebra.  The two sets of
quantities inhabit different rings because they are Mellin
transforms at different perturbative orders: $K$ at zeroth
order (free geometry), $D_p$ at $p$-th order (interacting geometry).

\paragraph{Summary.}
The Mellin transform \eqref{eq:mellin_bridge} is the silent engine
of the entire exchange constant taxonomy.  It converts discrete
spectral data (integer eigenvalues, integer degeneracies) into
transcendental content ($\pi$, $\zeta(s)$, Dirichlet $L$-values),
and the \emph{class} of transcendental that emerges is determined by
three properties of the operator whose heat kernel is being
transformed: (1)~whether it is first- or second-order (controlling
the motivic weight of the Mellin output); (2)~whether it acts on
scalars or spinors (modulating rational coefficients); and
(3)~whether the potential is free or interaction-deformed (controlling
whether the output stays in the $\pi$-polynomial ring or accesses the
full MZV algebra).  These are precisely the three axes of the QED
transcendental classification formalized in
Sec.~\ref{sec:taxonomy}.


%% ========================================================================
%% IV. TAXONOMY
%% ========================================================================

\section{Taxonomy of Exchange Constants}
\label{sec:taxonomy}

The following taxonomy classifies exchange constants by the type of
compactness mismatch they mediate.  The hierarchy---intrinsic,
calibration, embedding, flow, composition---can be read as increasing
departure from compactness: intrinsic constants convert between
compact manifolds of different dimension; calibration constants convert
between a compact manifold and flat space; embedding constants arise
from coordinate singularities in the non-compact description;
flow functions encode parameter-dependent symmetry breaking; and
composition constants measure the cost of factorizing a single compact
space into multiple incompatible coordinate systems.

Based on the catalog and the Weyl--Selberg identification, we
propose the following taxonomy:

\begin{table}[h]
\caption{Spectral-geometric exchange constants in the GeoVac
  hierarchy.  ``Source'' indicates the discrete structure;
  ``Target'' indicates the continuous coordinate system.
  ``Type'' classifies the exchange constant by what determines it.
  ``Algebraic (implicit)'' indicates that the function satisfies a
  polynomial equation $P(R,\mu) = 0$ with coefficients inherited from
  upstream exchange constants (Sec.~\ref{sec:algebraic_curve}).
  ``Piecewise'' indicates that the angular Hamiltonian has
  piecewise-smooth parameter dependence (from the split-region
  Legendre expansion), preventing a global characteristic polynomial.
  The function is piecewise-algebraic within each region.}
\label{tab:taxonomy}
\begin{ruledtabular}
\begin{tabular}{ccccc}
  Level & Source & Target & Constant & Type \\
  \hline
  --- & $S^1$ lattice & $\mathbb{R}^1$ & $\pi$ & intrinsic \\
  1 & $S^3$ graph & $\mathbb{R}^3$ & $\kappa = -1/16$ & conformal \\
  2 ($m{=}0$) & Laguerre rec. & prolate sph. & (none) & algebraic \\
  2 ($m{\neq}0$) & assoc. Laguerre & prolate sph. & $e^a E_1(a)$
    & embedding \\
  3 (FCI) & $S^3 {\times} S^3$ & (none) & (none) & algebraic \\
  3 (hyper.) & $\mathrm{SO}(6)$ reps & $S^5(R)$ param.
    & $\mu(R)$ & alg.\ (implicit) \\
  4 (prolate CI) & prolate product & (none) & (none) & algebraic \\
  4 (hyper.) & $\mathrm{SO}(6)$ reps & $(\rho,R)$ param.
    & $\mu(\rho,R)$ & piecewise \\
  5 (composed) & Level 3+4 fibers & full $N$-electron
    & PK pseudopot. & composition \\
  5 (balanced) & hydrogenic $\times$ multipole & cross-center $V_{ne}$
    & $\gamma(k,\alpha R)$ & algebraic \\
\end{tabular}
\end{ruledtabular}
\end{table}

The taxonomy reveals a hierarchy of exchange types:
\begin{enumerate}
  \item \textbf{Algebraic} (no exchange constant): When the
    computation stays on the graph or uses a basis whose inner
    product matches the graph structure, no transcendental
    appears.  Examples: Level~1 on $S^3$, Level~2 $\sigma$
    states, Level~3 FCI.

  \item \textbf{Intrinsic constants} ($\pi^{d/2}$): Determined
    by the manifold's dimension and volume alone, independent of
    any discretization or basis choice.  These are the Weyl
    constants---universal volume-to-eigenvalue conversion factors.
    Example: $\pi$ itself, for $S^1 \to \mathbb{R}$.

  \item \textbf{Conformal constants}: Determined by the conformal
    geometry of the projection from the compact graph manifold to
    flat space.  Example:
    $\kappa = -1/16 = -1/\Omega^4(0)$, where $\Omega(0) = 2$ is
    the stereographic conformal factor at $p = 0$
    (Sec.~\ref{sec:kappa_derivation}).  Equivalently derivable
    from graph vertex degree $d_{\max} = 4$; both routes yield
    the same value because the graph discretizes $S^3$.

  \item \textbf{Embedding constants}: Single transcendental
    numbers seeding algebraic recurrences.  Determined by the
    interaction between a spectral basis and a coordinate
    singularity.  Examples: $e^a E_1(a)$ for the prolate
    centrifugal pole; the $1/r_{12}$ electron-electron cusp,
    which is the embedding exchange constant for the two-electron
    coalescence in any product basis.
    The cusp embedding constant is unique in being \emph{removable}
    by a similarity transformation: the transcorrelated (TC)
    Jastrow factor $J = \tfrac{1}{2} r_{12}$ satisfies
    $e^{-J} (1/r_{12})\, e^{J} = \text{smooth}$, converting
    the algebraic partial-wave convergence $1/(l+\tfrac{1}{2})^4$
    to exponential.  This makes $1/r_{12}$ a ``soft'' embedding
    constant---present in any product-space representation but
    eliminable by a non-unitary change of basis---in contrast to
    $e^a E_1(a)$, which is structurally irreducible.

    The $1/r_{12}$ classification is \emph{universal}: the cusp
    condition survives relativistic corrections with only an
    $O((Z\alpha)^2)$ amplitude rescaling (Kutzelnigg-Morgan 1992), and
    partial-wave convergence exponents ($l^{-4}$ singlet, $l^{-6}$
    triplet) are identical in Schr\"odinger-Coulomb and Dirac-Coulomb
    (verified numerically at $Z=4$).  The Dirac program does not
    regularize the e-e cusp; full QED would be required.

    A \emph{distributional} sub-category appears in the Breit spin-spin
    interaction: the bare $1/r_{12}^3$ kernel diverges logarithmically
    at coalescence for every partial wave.  The Breit-Pauli tensor
    projection regularizes this, leaving $K_l^{BP} = r_<^l / r_>^{l+3}$
    with exact rational matrix elements for same-$n$ orbitals
    ($R_{BP}^2(2p,2p;2p,2p) = 7/768$ at $Z=1$).  For cross-$n$
    orbitals a new transcendental appears:
    $M_{\rm ret}^2(1s,2p;1s,2p) = 21344/729 - (10240/243)\log(2)$,
    the first $\log$-embedding constant in the framework.

    On the $S^3$ pair-state lattice (v2.9.1), 92--94\% of $V_{ee}$'s
    Frobenius mass is diagonal in the graph eigenbasis, with the
    off-diagonal residual concentrated on the $(1s,1s)$ coalescence
    pair-state.  The graph nearly diagonalizes $V_{ee}$ without
    containing $1/r_{12}$---direct evidence that the cusp is embedding
    content the graph cannot absorb, explaining the graph-native CI
    convergence floor (Paper~13, Track~DI).

  \item \textbf{Algebraic (implicit) functions}: Functions of
    continuous parameters defined implicitly by $P(R,\mu) = 0$,
    where $P$ is the characteristic polynomial of the
    parameter-dependent Hamiltonian $H(R) = H_0 + R \cdot V_C$.
    The coefficients of $P$ lie in $\mathbb{Q}(\pi, \sqrt{2})$---a
    transcendental extension of $\mathbb{Q}$---but the $R$-dependence
    itself is algebraic: $\mu(R)$ is a root of a polynomial in $R$.
    The transcendental content is \emph{inherited} from the angular
    coupling matrix $V_C$ (which contains upstream exchange constants
    of embedding type), not generated by the parameterization.
    Example: $\mu(R)$ for the $\mathrm{SO}(6) \to \mathrm{SO}(4)$
    symmetry crossover on $S^5$ at Level~3.  At $l_{\max} = 0$,
    $\mu(R)$ is an explicit linear function of $R$; at higher
    $l_{\max}$, it is an algebraic function of degree $l_{\max}+1$,
    with branch points at eigenvalue crossings in the complex
    $R$-plane.  At Level~4, the split-region Legendre expansion of
    the nuclear charge function introduces piecewise-smooth
    $\rho$-dependence at the region boundary $\rho = 1/2$, breaking
    the linear pencil structure. The angular eigenvalues $\mu(\rho)$
    are piecewise-algebraic within each region but do not satisfy a
    single global characteristic polynomial. The product-space CI
    saturation at 92.5\% $D_e$ remains the key irreducibility result.

  \item \textbf{Composition constants}: The accuracy cost of
    factorizing a full $N$-electron wavefunction into composed
    coordinates.  Unlike the preceding types, which convert between
    discrete and continuous descriptions, composition constants
    convert between a single-geometry and a multi-geometry
    (factorized) description.  The factorization destroys inter-group
    antisymmetry---the exchange operator $P_{ij}$ mixes coordinate
    systems that composition separates---and the PK pseudopotential
    is the local approximation to this nonlocal effect.  Example:
    PK for LiH composed geometry (Level~3 core + Level~4 valence),
    which achieves 144$\times$ angular compression at the cost of
    ${\sim}5\%$ $R_{\rm eq}$ error at $l_{\max} = 2$.
\end{enumerate}

The dependency hierarchy is explicit: intrinsic constants depend
only on dimension; calibration constants depend on the graph;
embedding constants depend on the basis; flow functions depend
on the full parameterized geometry; composition constants depend
on the factorization topology (which groups share coordinates).

\paragraph{Odd-zeta content from first-order spectral operators.}
The Dirac operator on unit $S^3$ (Camporesi--Higuchi~\cite{camporesi_higuchi1996}
spectrum $|\lambda_n| = n + 3/2$, degeneracy
$g_n^{\text{Dirac}} = 2(n+1)(n+2)$) introduces a structurally new
transcendental: the Fock Dirichlet series at the packing exponent
$s = d_{\max} = 4$ evaluates to
\begin{equation}\label{eq:dirac_dirichlet_zeta3}
  D_{g^{\text{Dirac}}}(4)
  \;=\; \sum_{m \geq 1} \frac{2m(m+1)}{m^4}
  \;=\; 2\,\zeta_R(2) + 2\,\zeta_R(3)
  \;=\; \frac{\pi^2}{3} + 2\,\zeta_R(3),
\end{equation}
a sympy-exact identity.  The two summands reflect the two homogeneity
classes of $g_m^{\text{Dirac}} = 2m^2 + 2m$: the weight-$m^2$ subchannel
produces $2\zeta_R(2)$ (the scalar $F$ content), while the weight-$m$
subchannel produces $2\zeta_R(3)$---a channel absent from the scalar
Fock degeneracy $g_n = n^2$.

This odd-zeta content is categorically different from calibration-tier
content.  Calibration constants carry even-zeta or $\pi^{\text{even}}$
content arising from second-order Riemannian operators through
Jacobi-$\theta$ inversion; the odd-zeta $\zeta_R(3)$ enters through
the first-order Dirac operator via half-integer Hurwitz-zeta identities,
and is $\mathbb{Q}$-linearly independent of $\zeta_R(2)$ by Ap\'ery's
theorem.  Within Paper~2's $K = \pi(B + F - \Delta)$, the odd-zeta tier
is a byproduct rather than a summand: $\zeta_R(3)$ is structurally
present on $S^3$ without participating in $K$.

\paragraph{Spinor-intrinsic content from first-order operators on a
spinor bundle.}
The Breit--Pauli reduction of the Dirac equation introduces a second
new transcendental tier: $\alpha^2$ and $\gamma := \sqrt{1 - (Z\alpha)^2}$.
The leading spin-orbit coupling
\begin{equation}\label{eq:bp_so}
  H_{\mathrm{SO}}
  \;=\; \tfrac{Z\alpha^2}{2}\,\tfrac{1}{r^3}\,\mathbf{L}\cdot\mathbf{S},
\end{equation}
evaluated in the $(\kappa, m_j)$ basis via Szmytkowski's reduction
$\langle \mathbf{L}\cdot\mathbf{S}\rangle = -(\kappa+1)/2$ and the
hydrogenic $\langle 1/r^3\rangle_{n,l}
= Z^3/[n^3\, l(l+\tfrac{1}{2})(l+1)]$, gives
\begin{equation}\label{eq:bp_so_matrix}
  \langle n,\kappa,m_j|\,H_{\mathrm{SO}}\,|n,\kappa,m_j\rangle
  \;=\; -\,\frac{Z^4\,\alpha^2\,(\kappa+1)}{4\,n^3\, l(l+\tfrac{1}{2})(l+1)},
\end{equation}
with Kramers cancellation at $l=0$.  Full-Dirac radial corrections enter
through $\gamma$, which satisfies
\begin{equation}\label{eq:gamma_def}
  \gamma^2 + (Z\alpha)^2 \;=\; 1,
\end{equation}
making $\gamma$ algebraic of degree~2 over $\mathbb{Q}(\alpha^2)$.
Together $\alpha^2$ and $\gamma$ generate the ring
\begin{equation}\label{eq:ring_sp}
  R_{\mathrm{sp}} \;:=\; \mathbb{Q}\bigl(\alpha^2\bigr)
    \bigl[\gamma\bigr]/\bigl(\gamma^2 + (Z\alpha)^2 - 1\bigr),
\end{equation}
in which every spinor-intrinsic coefficient lies---mechanically
certified by \texttt{geovac/spinor\_certificate.py}.

This tier is categorically distinct from both the odd-zeta tier and
calibration: $\alpha^2$ is a \emph{coupling} content appearing in
Hamiltonian coefficients (not in the basis or spectrum), depending only
on the existence of spin coupling in the operator.  Higher-order
corrections (Darwin $\alpha^4$, mass-velocity $\alpha^4$, Breit
$\alpha^2$ in $V_{ee}$) populate the same ring $R_{\mathrm{sp}}$ without
creating new transcendental tiers.

\paragraph{Operator-order $\times$ bundle classification.}

\begin{table}[h]
\caption{Transcendental content classified by operator order and bundle
  type.  The (2nd-order, spinor) cell is degenerate with scalar
  calibration ($\pi^{\mathrm{even}}$ only)---see the $D^2$ discussion
  below.}
\label{tab:bundle_operator_tiers}
\begin{ruledtabular}
\begin{tabular}{llll}
  Tier & Example transcendentals & Operator order & Bundle \\
  \hline
  Intrinsic         & $\pi^{d/2}$ Weyl prefactors & --- & scalar \\
  Calibration       & $\kappa = -1/16$, $F = \pi^2/6$
                    & 2nd ($\Delta_{\text{LB}}$)
                    & scalar \\
  Embedding         & $e^a E_1(a)$, $1/r_{12}$
                    & ---
                    & graph $\to$ $\mathbb{R}^3$ \\
  Flow / composition & $\mu(\rho,R)$, PK
                    & ---
                    & multi-geometry \\
  Odd-zeta          & $\zeta_R(3), \zeta_R(5), \dots$
                    & 1st ($D$)
                    & Dirac on curved space \\
  Spinor-intrinsic  & $\alpha^2$, $\gamma$
                    & 1st ($D$ / Breit--Pauli)
                    & spinor bundle \\
  $D^2$ calibration & $\pi^{\text{even}}$ (degenerate)
                    & 2nd ($D^2$)
                    & spinor bundle \\
\end{tabular}
\end{ruledtabular}
\end{table}

\noindent Tier~1 ($\zeta$) enters through the \emph{spectrum} of the
Dirac operator; Tier~2 ($\alpha^2$, $\gamma$) enters through the
\emph{coupling} of spin to orbital angular momentum in the Hamiltonian
coefficients.  Neither introduces $\pi$ beyond the calibration tier.

\paragraph{The squared Dirac operator and the operator-order discriminant.}
The squared Dirac operator $D^2$ on unit $S^3$ fills the last open cell
of Table~\ref{tab:bundle_operator_tiers}.  By the Lichnerowicz
formula~\cite{lichnerowicz1963}, $D^2 = \nabla^*\nabla + R/4$ with
$R/4 = 3/2$ (rational shift).  Its eigenvalues $(n+3/2)^2$ yield a
spectral zeta function with an exact closed form at every integer~$s$:
\begin{equation}\label{eq:zeta_d2}
  \zeta_{D^2}(s)
  \;=\; 2^{2s-1}\bigl[\lambda(2s{-}2) - \lambda(2s)\bigr],
\end{equation}
where $\lambda(a) = (1 - 2^{-a})\,\zeta_R(a)$ is the Dirichlet lambda
function.  Since $\zeta_R(2k) = \text{rational} \times \pi^{2k}$
(Bernoulli), every term is a polynomial in $\pi^2$ with rational
coefficients.  \emph{No odd-zeta content appears at any~$s$}: squaring
maps odd-power spectral sums to even-power sums, structurally
eliminating the $\zeta_R(3)$ found in the first-order Dirac zeta
(Eq.~\ref{eq:dirac_dirichlet_zeta3}).

The (2nd-order, spinor) cell is therefore \textbf{degenerate with scalar
calibration}: both produce $\pi^{\text{even}}$ exclusively.  Operator
order---first vs.\ second---is the primary transcendental discriminant
on round $S^3$.  Bundle type modulates degeneracies, eigenvalue lattices,
and rational coefficients, but not the transcendental class.  The
taxonomic grid collapses to three effective tiers:
\begin{enumerate}
  \item \textbf{First-order operators:} odd-zeta content
    ($\zeta_R(3), \zeta_R(5), \ldots$) plus spinor coupling content
    ($\alpha^2$, $\gamma$).
  \item \textbf{Second-order operators:} $\pi^{\text{even}}$ exclusively,
    on both scalar and spinor bundles.
  \item \textbf{Basis and composition tiers:} embedding, flow,
    composition constants (operator-independent).
\end{enumerate}
This is consistent with Paper~24's Coulomb/HO asymmetry: second-order
Riemannian operators introduce $\pi$ through nonlinear conformal
projections, while first-order operators (the HO Bargmann Euler operator
on $H^2(S^5)$, the Dirac operator on $S^3$) have $\pi$-free spectra
in rational arithmetic.

\paragraph{Motivic weight and operator order.}
The operator-order discriminant---second-order operators produce
$\pi^{\text{even}}$ content, first-order operators produce odd-zeta
content---has a natural interpretation in the theory of mixed Tate
motives over $\mathrm{Spec}(\mathbb{Z})$.  In this framework
(Goncharov~\cite{goncharov1999}, Zagier~\cite{zagier1994}), the Riemann
zeta value $\zeta_R(n)$ is a \emph{period} of the mixed Tate motive
$\mathfrak{M}(n) = \mathbb{Q}(-n)$ of weight~$n$.  Even zeta values
$\zeta_R(2k) = \text{rational} \times \pi^{2k}$ are periods of motives
that factor through the de~Rham cohomology of even-dimensional compact
manifolds (the Bernoulli formula).  Odd zeta values $\zeta_R(2k{+}1)$
are periods of genuinely independent motives: they do not reduce to
powers of~$\pi$, and $\zeta_R(3)$ in particular is irrational (Ap\'ery
1978) and conjectured algebraically independent from~$\pi$.

The GeoVac operator-order discriminant maps onto this motivic
partition:
\begin{center}
\begin{tabular}{llll}
  \hline
  Operator order & Motivic weight & Transcendental class & Example \\
  \hline
  2nd ($\Delta_{\text{LB}},\; D^2$)
    & even ($2k$) & $\zeta_R(2k) = \text{rat} \cdot \pi^{2k}$
    & $\kappa$, $\mathrm{Vol}(S^3)$, $a_2$ \\
  1st ($D$)
    & odd ($2k{+}1$) & $\zeta_R(2k{+}1) \perp \mathbb{Q}(\pi)$
    & D3: $\zeta_R(3)$ in Dirac Dirichlet \\
  \hline
\end{tabular}
\end{center}

At motivic weight~3, $\zeta_R(3)$ appears as a period of
$\mathbb{Q}(-3)$ in multiple independent geometric contexts:
the Dirac Dirichlet series on~$S^3$ (this work,
Eq.~\ref{eq:dirac_dirichlet_zeta3}), hyperbolic 3-manifold volumes
(Goncharov~\cite{goncharov1999}), two-loop QED $g{-}2$
(Petermann~\cite{petermann1957}; Sommerfield~\cite{sommerfield1957}),
and Chern--Simons theory / algebraic K-theory
($K_5(\mathbb{Z})$ regulator).  The common thread is~$S^3$ or its
quotients.  Zagier's dimension conjecture~\cite{zagier1994} gives
$d_3 = 1$: $\zeta_R(3)$ is the unique independent transcendental
at weight~3, consistent with its single appearance in the GeoVac
Dirac sector.

The connection to QED loop order is structurally predicted by the
motivic correspondence.  The T9 theorem
(Eq.~\ref{eq:zeta_d2}) guarantees that $\zeta_R(3)$ cannot enter through
any single trace of a squared propagator $D^{-2s}$---i.e., at one loop.
It enters at two loops, where irreducible products of first-order
propagators $D^{-1} \otimes D^{-1}$ produce odd motivic weight.  The
operator-order discriminant thus determines the \emph{loop order} at
which new transcendental content first appears: one loop accesses only
even-weight motives ($\pi^{2k}$); two loops access odd-weight motives
($\zeta_R(3)$, $\zeta_R(5)$, \ldots).

The even/odd Hurwitz spectral zeta of the first-order Dirac operator
provides a sharp computational check:
\begin{equation}\label{eq:hurwitz_discriminant}
  D(s) \;=\; 2\bigl(2^{s-2}-1\bigr)\,\zeta_R(s{-}2)
           - \tfrac{2^s - 1}{2}\,\zeta_R(s).
\end{equation}
This exact formula separates the even and odd content at each~$s$:
at $s = 4$, $D(4) = 2\zeta_R(2) + 2\zeta_R(3)$, reproducing
Eq.~(\ref{eq:dirac_dirichlet_zeta3}).  At even~$s$, both terms are
$\pi^{\text{even}}$; at odd~$s$, the $\zeta_R(s{-}2)$ and
$\zeta_R(s)$ terms are independently transcendental.  The QED vacuum
polarization at one loop involves the Seeley--DeWitt coefficients
$a_0$, $a_1$, $a_2$ of~$D^2$, which by
Eq.~(\ref{eq:zeta_d2}) are purely $\pi^{\text{even}}$; the
Schwinger~\cite{schwinger1948} result $a_e^{(1)} = \alpha/(2\pi)$
is a single power of~$\pi$, consistent with one-loop = even motivic
weight.  At two loops, $\zeta_R(3)$ appears in the electron $g{-}2$
(Petermann~\cite{petermann1957}; Sommerfield~\cite{sommerfield1957}), consistent with
the discriminant's prediction that odd-weight content requires
irreducible products of first-order propagators.

This motivic interpretation refines the compactness thesis of
Sec.~\ref{sec:intro}: exchange constants are not only projections from
compact to non-compact geometry---they are \emph{periods} of specific
mixed Tate motives, and the operator order on the compact space
determines the motivic weight of the resulting transcendental content.

\paragraph{Dirichlet $L$-values from vertex topology.}
The operator-order discriminant classifies free-propagator content.
A third axis emerges when interaction vertices are included.  The
QED vertex on~$S^3$ couples Dirac spinor harmonics at levels $n_1$,
$n_2$ to a transverse vector harmonic at level~$n_\gamma$ with
SO(4) selection rules requiring $n_1 + n_2 + n_\gamma$ odd (the
$\gamma^\mu$ coupling flips parity).  When this selection rule is
imposed on the two-loop sunset diagram, it splits the Dirac Dirichlet
series into even- and odd-$n$ sub-sums with structurally different
transcendental content.  At $s = 4$ (the packing exponent), the
exact Hurwitz decomposition gives:
\begin{align}
  D_{\text{even}}(4) &= \tfrac{\pi^2}{2} - \tfrac{\pi^4}{24}
    - 4G + 4\beta(4), \label{eq:d_even4} \\
  D_{\text{odd}}(4) &= \tfrac{\pi^2}{2} - \tfrac{\pi^4}{24}
    + 4G - 4\beta(4), \label{eq:d_odd4}
\end{align}
where $G = \sum_{n \ge 0} (-1)^n/(2n{+}1)^2 \approx 0.9159$ is
Catalan's constant and $\beta(4) = \sum_{n \ge 0} (-1)^n/(2n{+}1)^4
\approx 0.9889$ is the Dirichlet beta function at~4.  Both are values
of the $L$-function $L(s, \chi_{-4})$ for the unique non-principal
character mod~4.  Crucially, $G$ and $\beta(4)$ cancel in the full sum
$D_{\text{even}} + D_{\text{odd}} = D(4) = \pi^2 - \pi^4/12$---the
Dirichlet content is invisible without vertex restriction.  The
mechanism is the Hurwitz quarter-shift: even-$n$ modes involve
$\zeta(s, 3/4)$ and odd-$n$ modes $\zeta(s, 5/4)$, whose difference
is $4^s \beta(s)$ (a Dirichlet $L$-value), while their sum reduces
to Riemann zeta ($\pi^{\text{even}}$).

This extends the taxonomy to three axes:
(1)~\emph{operator order} (1st~vs 2nd) determines motivic weight
(odd~vs even~zeta);
(2)~\emph{bundle type} (scalar~vs spinor) modulates coefficients
within the same weight class;
(3)~\emph{diagram topology}---specifically, vertex parity from the
interaction structure---introduces Dirichlet $L$-function characters.
The Dirichlet $\beta$-values are a genuinely new transcendental class:
they arise from the \emph{interaction} (how electron and photon modes
couple at the vertex), not from the free propagators.  In the GeoVac
exchange-constant taxonomy, they represent a ``vertex exchange
constant'' tier that lives beyond both calibration~$\pi$ and
Riemann odd-zeta.  Verified by PSLQ at 80-digit precision;
35/35 tests pass in \texttt{tests/test\_qed\_vertex.py}.

\paragraph{Photon vector structure as a calibration exchange constant.}
The graph-native QED construction (Paper~28, GN-1 through GN-7)
computes one-loop QED entirely on the finite Fock graph using exact
algebraic arithmetic.  On the graph, the photon lives on edges
(1-cochains) as a \emph{scalar} object---the edge Laplacian
$L_1 = B^T\!B$ propagates it without reference to any vector
structure.  In the continuum, the physical photon is a transverse
vector field on~$S^3$ carrying $\mathrm{SO}(4) = \mathrm{SU}(2)_L
\times \mathrm{SU}(2)_R$ quantum numbers, and this vector structure
enforces selection rules (vertex parity, SO(4) channel counting~$W$)
that are absent on the graph.  Adding explicit photon angular momentum
labels $(q, m_q)$ with Clebsch--Gordan vertex coupling and E1 parity
enforcement recovers 7~of~8 continuum QED selection rules for
\emph{scalar} electrons (Paper~28).  The sole failure for scalar
electrons is Furry's theorem: diagonal vertex couplings
$V(a,a,q,m_q) \neq 0$ for $l \geq 1$ when $q$ is odd.  For
\emph{Dirac} electrons, Furry's theorem is recovered as well,
achieving 8/8 (see below).

The calibration cost of the $1 \to 7$ transition is precisely one
transcendental per loop: $\pi$, entering as the $S^2$ Weyl exchange
constant $1/(4\pi)$ from the vector spherical harmonic normalization
$\sqrt{(2l+1)(2q+1)(2l'+1)/(4\pi)}$.  This is a concrete instance
of the calibration exchange constants cataloged in the taxonomy table
(Table~\ref{tab:taxonomy}): the $S^2$ Weyl constant $\omega_2/(2\pi)^2
= 1/(4\pi)$ converts the discrete angular momentum quantum numbers
into the continuous solid angle measure required for vector photon
normalization.  The entire scalar graph QED is $\pi$-free
($\mathbb{Q}[\sqrt{2}, \sqrt{3}, \sqrt{6}, \ldots]$); the vector
photon introduces exactly the $\omega_2 = \pi$ entry from the
Weyl constant table.

In the taxonomy, the photon's angular momentum structure therefore
sits in the \textbf{calibration} tier, alongside the volume
normalization $\mathrm{Vol}(S^3) = 2\pi^2$ and the Hurwitz-zeta
regularization of infinite spectral sums: it is a property of the
\emph{manifold embedding}, not of the graph.  This sharpens the
graph/continuum partition for QED: the graph captures the
algebraic/combinatorial content (rational propagators, algebraic CG
vertices, $\pi$-free observables), while the continuum projection
adds three calibration-tier ingredients: (1)~volume normalization,
(2)~infinite-sum regularization, and (3)~photon vector structure
(cost: $1/(4\pi)$ per loop).
The 8~continuum selection rules partition by transcendental cost
into three tiers.  For \emph{scalar} electrons, the partition is
$1$--$7$--$8$: graph topology (1~rule, zero cost)
$\to$ angular momentum conservation (6~rules, calibration cost
$1/(4\pi)$) $\to$ field theory (1~rule, requires second-quantized
$C$-symmetry).  For \emph{Dirac} electrons, the partition refines
to $1$--$6$--$1$--$8$: graph topology (1~rule, intrinsic, zero cost)
$\to$ angular momentum conservation (6~rules, calibration $1/(4\pi)$)
$\to$ Dirac spinor phase constraint (1~rule, intrinsic, zero cost).

\paragraph{The Dirac spinor phase constraint as an intrinsic tier.}
When the electron is described by a Dirac spinor
$\psi = (f \cdot \Omega_\kappa,\; i g \cdot \Omega_{-\kappa})^T$
(Paper~23), the $\alpha$-matrix vertex coupling is off-diagonal in the
large/small component space.  The factor of~$i$ on the small component
makes the diagonal coupling $V(a,a,q,m_q)$ vanish identically for all
states: $V_{LS} + V_{SL} = 0$ by Hermiticity.  This is the mechanism
by which Furry's theorem is recovered---tadpole diagrams (self-loops at
vertex~$a$) vanish because the single-particle Dirac vertex is
kinematically off-diagonal.

The key taxonomic point is the \emph{cost} of this recovery.  The
$1 \to 7$ transition (adding photon angular momentum labels) costs
exactly $\pi$ per loop---it is calibration content.  The $7 \to 8$
transition (Dirac spinor phase) costs \emph{nothing}: no new
transcendental enters, no continuum embedding is required.  The spinor
phase constraint is a kinematic identity of the single-particle Dirac
vertex, not a field-theoretic symmetry ($C$-conjugation).  It lives in
the same \textbf{intrinsic} tier as the graph topology rule (Gaunt/CG
sparsity), not in the calibration tier.

This distinction has a concrete implication: the $7 \to 8$ transition
does \emph{not} require second-quantized field theory.  Any
single-particle framework using the Dirac equation on~$S^3$ with
vector-photon coupling automatically recovers all~8 selection rules.
The standard attribution of Furry's theorem to charge conjugation
symmetry is operationally correct in flat-space QFT but obscures the
fact that, on the compact $S^3$ graph, the same selection rule has a
simpler origin in the Dirac spinor structure itself.

\paragraph{Flat-space consistency of the transcendental classification.}
The three-axis taxonomy is derived from $S^3$ spectral sums at finite
radius.  A natural concern is whether the classification is a
finite-curvature artifact.  We verify that it is not.
On $S^3$ of radius~$R$, the Dirac eigenvalues scale as
$|\lambda_n(R)| = (n + 3/2)/R$ while the degeneracies $g_n$ are
$R$-independent.  The Dirichlet series therefore satisfies
$D(s,R) = R^s\,D(s,1)$ \emph{exactly}, making the ratio
$D(s,R)/R^s$ an $R$-independent structural invariant.  In the
$R \to \infty$ limit, Weyl's law confirms that the mode count
$N(n_{\max}) = 2(n_{\max}+1)(n_{\max}+2)(n_{\max}+3)/3$
matches the flat-space momentum-integral density
$(2/3)\,R^3 p_{\max}^3$, verifying the spectral-to-momentum
correspondence.  The even/odd $s$-parity discriminant persists at
every~$R$: $D(s_{\text{even}},R)/R^s$ is $\pi^{\text{even}}$;
$D(s_{\text{odd}},R)/R^s$ contains odd-zeta.  The \emph{same}
transcendental constants---$\zeta_R(3)$, $G$, $\beta(4)$---that
appear in the $S^3$ spectral sums also appear in flat-space
two-loop QED (Petermann~\cite{petermann1957};
Sommerfield~\cite{sommerfield1957}), confirming that curvature
modifies rational coefficients but not the transcendental class.
An honest limitation: matching the full Petermann--Sommerfield
coefficient $a_2 = -197/144 + \pi^2/12 + 3\zeta_R(3)/4
- \pi^2 \ln 2 / 2$ requires diagram-specific computation
(renormalization subtractions, SO(4) Clebsch--Gordan algebra)
beyond the scope of structural spectral sums.
The flat-space limit thus validates the taxonomy as a
\emph{structural} property of the operator algebra, independent
of the $S^3$ compactification scale.

\paragraph{Irreducibility of the Level~4 piecewise structure.}
The piecewise-smooth $\rho$-dependence of $\mu(\rho,R)$ at Level~4
is not a coordinate artifact: it cannot be resolved by geometric
elevation to any higher-dimensional space (Track~Z).  Three approaches
were ruled out: (i)~algebraic blow-up of $|s - \rho|$ produces smooth
sheets but with $\rho$-dependent integration limits ($\alpha = \arccos\rho$,
$\arcsin\rho$), making matrix elements transcendental in~$\rho$;
(ii)~Lie algebra elevation fails because $\min(s,\rho)/\max(s,\rho)$
is $C^0$ but not~$C^1$, incompatible with matrix elements of any
finite-dimensional Lie group representation (which are real-analytic);
(iii)~$S^3 \times S^3$ product space fails because two-center nuclear
attraction cannot be smooth on any single compact manifold.
The piecewise structure thus belongs to the embedding category
in the taxonomy: it is the irreducible geometric content of placing
two nuclei at finite separation in hyperspherical coordinates, analogous
to the $1/r_{12}$ codimension obstruction (Track~W) but affecting the
nuclear rather than the electronic coupling.

\paragraph{Even-odd staircase and channel convergence rate.}
An extended convergence study at Level~4 (Paper~15, Track~AA)
reveals a pronounced even-odd staircase in $D_e$ vs $l_{\max}$:
even increments produce large gains ($+42$, $+6.8$, $+1.7$~pp)
while odd increments produce near-zero gains ($+0.66$, $+0.08$~pp).
The root cause (Track~AC) connects directly to the exchange constant
structure: the nuclear coupling matrix
$\langle l_1' l_2' | V_{\rm nuc} | l_1 l_2 \rangle$ is diagonal
in one electron's angular momentum index (Kronecker delta), so
direct coupling from the ground channel $(0,0)$ requires a
$\Delta l = 2$ transition enforced by the gerade constraint
$l_1 + l_2 = \text{even}$.  Odd-$l_{\max}$ channels couple only
indirectly through the electron-electron multipole expansion.

This selection rule constrains the \emph{convergence rate} of the
partial-wave expansion: the nuclear coupling---a piecewise exchange
constant in the taxonomy---couples to the ground state only through
quadrupolar ($\Delta l = 2$) transitions, producing an algebraic
convergence rate $\sim (l + 1/2)^{-2}$ with an even-odd modulation.
The $\sigma$ and $\pi$ channel sectors are completely decoupled
(both nuclear and $e$--$e$ coupling conserve $m$), with $\pi$
contributing a constant ${\sim}6.6$~pp additive offset.

\paragraph{PK pseudopotential as composition exchange constant.}
The Phillips--Kleinman pseudopotential in the composed geometry
(Paper~17) enforces core-valence orthogonality, which is a
Pauli exclusion effect---distinct from the Coulomb exchange
constants cataloged above.  Without PK, no equilibrium
geometry exists (the PES is monotonically attractive), establishing
PK as essential for creating the bonding minimum.  Ab initio
derivation of the $R$-dependent PK scaling parameter $R_{\rm ref}$
from core properties fails due to a fundamental scale mismatch:
all core-derived candidates ($0.27$--$0.82$~bohr) are too small
by a factor of $4$--$11$.

A complete survey of all solver $\times$ PK combinations
(Paper~17, Table~VI) confirms that PK is not an improvable
approximation.  The 2D variational solver eliminates ${\sim}75\%$
of the $l_{\max}$ drift (from the adiabatic approximation), but
$l$-dependent PK has zero additional effect on the 2D solver.
The residual drift ($+0.15$~bohr/$l_{\max}$) is intrinsic to
the composed-geometry factorization.

Three graph-native alternatives to PK were investigated and all
failed: (i)~node exclusion between the core $S^3$ and valence
Level~4 graphs has no coordinate correspondence;
(ii)~the fiber bundle Berry connection is intra-group, and the
exchange operator $P_{ij}$ mixes coordinate systems;
(iii)~spectral orthogonality (function-space) does not imply
coordinate-space exclusion.  The master obstruction is that
antisymmetry requires a shared coordinate system; composition
factorizes $\Psi_{\rm total}$ into incompatible coordinates.

PK therefore defines a sixth exchange type in the taxonomy:
\textbf{composition constants} convert between a single-geometry
(full $N$-electron) and a composed-geometry (factorized) description.
Unlike the other exchange constants, which convert between discrete
and continuous descriptions, composition constants measure the
accuracy cost of the factorization itself.  The ``exchange rate''
is quantified by the 144$\times$ angular compression that composed
LiH achieves over a full 4-electron solver on $S^{11}$
(Paper~17, Sec.~V): the ${\sim}5\%$ $R_{\rm eq}$ error at
$l_{\max} = 2$ is the irreducible cost of this compression.

This classification is validated by the full 4-electron Level~4N
solver (Paper~17, Sec.~VIII): exact $S_4$ antisymmetry produces
molecular equilibrium without PK at $l_{\max} \geq 2$, confirming
that PK approximates---rather than creates---the physics of Pauli
exclusion.  The composition exchange constant is the cost of
replacing exact inter-group antisymmetry with a pseudopotential
in the factorized coordinate system.


%% ========================================================================
%% V. THE FINE STRUCTURE CONSTANT CONNECTION
%% ========================================================================

\section{Connection to the Fine Structure Constant}
\label{sec:alpha}

Paper~2~\cite{loutey_paper2} observes that the fine structure
constant $\alpha$ satisfies the cubic $\alpha^3 - K\alpha + 1 = 0$,
where the coefficient
\begin{equation}
  K = \pi\!\left(42 + \frac{\pi^2}{6} - \frac{1}{40}\right)
  = 137.036064
  \label{eq:K}
\end{equation}
is constructed from spectral invariants of the Hopf fibration
$S^1 \to S^3 \to S^2$.  The combination rule
$K = \pi(B + F - \Delta)$ was identified as conjectural---an
empirical formula with structural support but no first-principles
derivation.

We observe that this combination rule has precisely the structure
of a spectral-geometric exchange:
\begin{itemize}
  \item $B = 42$ is discrete representation-theoretic data:
    the degeneracy-weighted $\mathrm{SO}(3)$ Casimir trace over
    the first three shells of $S^3$.
  \item $F = \zeta(2) = \pi^2/6$ is a continuous spectral
    invariant: the spectral zeta function of $S^1$ at $s = 1$.
  \item $\Delta = 1/40$ is a boundary correction mixing discrete
    eigenvalue data ($|\lambda_3| = 8$) with the cumulative state
    count ($N(2) = 5$).
  \item The overall factor $\pi$ is $\omega_2 = \pi$, the 2D
    ball volume, which serves as the numerator of the $S^2$ Weyl
    density and as the packing-plane prefactor
    $\sigma_0 = \pi d_0^2 / 2$ in Paper~0 (the base of the Hopf
    fibration, the photon propagation sphere).
\end{itemize}

The structure of $K$ is thus: the 2D ball volume $\omega_2 = \pi$
(the numerator of the $S^2$ Weyl density, and the same $\pi$ that
fixes Paper~0's packing-plane prefactor $\sigma_0 = \pi d_0^2 / 2$)
multiplying a sum of the $S^2$ Casimir trace ($B$, discrete),
the $S^1$ spectral zeta value ($F$, transcendental), and an
$S^3$ boundary term ($\Delta$, mixed).  This is a Weyl-type
formula that converts between the discrete bundle
(representation labels on $S^3$) and the continuous coupling
(the electromagnetic interaction strength measured in
$\mathbb{R}^3$).

If this identification is correct, then the combination rule
$K = \pi(B + F - \Delta)$ is not arbitrary numerology but an
instance of the same spectral-geometric exchange that produces
$\kappa = -1/16$ at Level~1 and $e^a E_1(a)$ at Level~2.
The exchange constant for the Hopf bundle is $K$ itself---the
``cost'' of projecting the $S^3$ electron state space through
the $S^1$ gauge fiber onto the $S^2$ photon base, measured
in units of coupling strength.

\begin{conjecture}
The combination rule $K = \pi(B + F - \Delta)$ is a
spectral-geometric exchange constant for the Hopf fibration,
analogous to the Weyl--Selberg constants for compact manifolds.
The specific form---$\pi$ multiplying a sum of Casimir data
and zeta values---should be derivable from the spectral geometry
of $S^1 \to S^3 \to S^2$, potentially via the heat kernel
expansion or the analytic torsion of the Hopf bundle.
\end{conjecture}

Two additional observations support this identification:

\begin{observation}[Structural match]
The combination rule $K = \pi(B + F - \Delta)$ mixes discrete
data ($B = 42$, from representation theory) with continuous
spectral data ($F = \zeta(2) = \pi^2/6$, from the $S^1$
spectral zeta function), multiplied by an overall geometric
factor ($\omega_2 = \pi$, the 2D ball volume that enters as the
numerator of the $S^2$ Weyl density and as the packing-plane
prefactor $\sigma_0 = \pi d_0^2 / 2$ in Paper~0, with $S^2$ the
Hopf base).  This is precisely the pattern identified throughout this
paper: a transcendental conversion factor mediating between
discrete and continuous descriptions of the same geometry.
The structural match is evidence---not proof---that the
combination rule is an instance of the exchange constant pattern
rather than numerology.
\end{observation}

\begin{observation}[$\alpha$ is likely transcendental]
The fine structure constant $\alpha$ is a root of the cubic
$\alpha^3 - K\alpha + 1 = 0$, where $K$ involves $\pi$ and
$\zeta(2) = \pi^2/6$.  Roots of polynomials with transcendental
coefficients \textit{can} be algebraic (e.g., $x^2 - \pi x +
\pi(\pi - 1) = 0$ has the algebraic root $x = 1$), but this
requires special cancellations that are generically absent.
Since $K$ is a non-trivial combination of $\pi$ and $\pi^2$
with rational coefficients, it is overwhelmingly likely
(though not proven) that $\alpha$ is transcendental---consistent
with its role as a spectral-geometric exchange constant.
\end{observation}

\subsection{Three-tier structure of the combination rule}
\label{sec:alpha_three_tier}

An ongoing computational investigation into the origin of the three
components of $K = \pi(B + F - \Delta)$---using exact-rational
symbolic computation and high-precision arithmetic---has pinned down
$B$ and $F$ as genuinely different objects on the Fock $(n,l)$
lattice, and has eliminated a broad family of candidate mechanisms
for $\Delta$. We observe that the three terms appear to occupy
three categorically distinct tiers of the exchange-constant
taxonomy rather than specializing a single common generator. We
summarize the four structural findings below; none of them alters
the conjectural status of Paper~2's cubic identity.

\paragraph{1. A $\kappa$--$B$ identity from the Fock weight.}
The Fock projection of the hydrogenic basis onto $S^3$ carries a
residual momentum-space weight
$w(p) = (p^2 + p_0^2)^{-2}$
after the overall $(2p_0)^{5/2}$ prefactor of the Fock map is
stripped.  Computation reveals that the diagonal matrix elements
$\langle n,l|w|n,l\rangle_{p_0 = 1}$ for $n \le 3$ are exact
rationals whose denominators divide $32 = 2\,d_{\max}^2$:
\begin{equation}
  \{\langle n,l|w|n,l\rangle_{p_0 = 1}\}_{(n,l)}
  = \Bigl\{ \tfrac{7}{16},\, \tfrac{5}{8},\, \tfrac{3}{8},\,
            \tfrac{5}{8},\, \tfrac{17}{32},\, \tfrac{11}{32} \Bigr\}
  \label{eq:fock_weight_values}
\end{equation}
for $(n,l) \in \{(1,0),(2,0),(2,1),(3,0),(3,1),(3,2)\}$.  Weighted
by the Hopf base factor $(2l+1)\,l(l+1)$ and summed,
\begin{equation}
  \sum_{n=1}^{3}\sum_{l=0}^{n-1} (2l+1)\,l(l+1)\,
    \langle n,l|w|n,l\rangle_{p_0 = 1}
  = \frac{63}{4}
  = 6\,B\,|\kappa|,
  \label{eq:kappa_B_identity}
\end{equation}
with $B = 42$ the Casimir trace of Paper~2 and $\kappa = -1/16$
the calibration constant of Sec.~\ref{sec:kappa_derivation}.
Equivalently, averaging over the six $(n,l)$ cells gives
$B\,|\kappa| = 21/8$.  This is the first exact rational algebraic
link in the GeoVac framework between the calibration constant and
the Hopf base invariant: $|\kappa| = 1/d_{\max}^2$ appears as the
universal denominator factor of the Fock-weight matrix elements,
while $B$ appears in the Hopf-weighted sum.  The identity
(\ref{eq:kappa_B_identity}) is verified by exact-rational symbolic
computation with no fitted parameters.

\paragraph{2. A Dirichlet identity for $F$.}
Define the Dirichlet series of the $S^3$ Fock shell degeneracy
$g_n = n^2$,
\begin{equation}
  D_{n^2}(s) \;:=\; \sum_{n=1}^{\infty} g_n\, n^{-s}
             \;=\; \sum_{n=1}^{\infty} n^{-(s-2)}
             \;=\; \zeta_R(s - 2),
  \label{eq:Dn2_def}
\end{equation}
evaluated at the packing exponent $s = d_{\max} = 4$ of Paper~0:
\begin{equation}
  D_{n^2}(d_{\max}) = \zeta_R(2) = \frac{\pi^2}{6} = F.
  \label{eq:F_dirichlet}
\end{equation}
The exponent $s = 4$ admits three independent natural readings:
$s = d_{\max}$ from Paper~0's packing axiom; $s = \dim(\mathbb{R}^4)$,
the ambient Euclidean space into which $S^3$ is stereographically
projected in Paper~7; and $s = 2 N_{\mathrm{init}}$ from the same
packing axiom.  The weight $n^2$ is the degeneracy of the $n$-th
$S^3$ shell in Fock's projection.  To our knowledge
(\ref{eq:F_dirichlet}) is the first identification of $F$ from a
graph-intrinsic quantity in the GeoVac framework---that is, $F$ is
produced here without invoking an embedding metric or a
sphere-spectral construction.  No other natural Dirichlet series
on the shell lattice (per-shell Casimir, state count, Laplacian
eigenvalue, $l$-max weight) produces an exact hit at any integer
$s$ in the range $[3, 12]$.

\paragraph{3. Three-tier classification.}
Combining the $\kappa$--$B$ identity (\ref{eq:kappa_B_identity})
with the Dirichlet identity (\ref{eq:F_dirichlet}) and the
Paper~2 canonical form $1/\Delta = |\lambda_{n_{\max}}|\, N(n_{\max}{-}1)$,
the three components of $K = \pi(B + F - \Delta)$ have
categorically different origins:

\begin{center}
\small
\renewcommand{\arraystretch}{1.2}
\begin{tabular}{@{}lccl@{}}
\hline
Component & Value & Type & Origin \\
\hline
$B$      & $42$        & rational, finite sum       & Casimir trace
   truncated at $n_{\max}=3$ \\
$F$      & $\pi^2/6$   & transcendental, infinite   & Fock-degeneracy
   Dirichlet at $s = d_{\max}$ \\
         &             & series                     & \\
$\Delta$ & $1/40$      & rational, boundary product & $|\lambda_{n_{\max}}|
   \cdot N(n_{\max}{-}1)$ \\
\hline
\end{tabular}
\end{center}

\noindent The three objects do not arise from a single common
generator.  $B$ is a finite sum over the first three shells;
$F$ is an infinite Dirichlet series with no meaningful
finite truncation at $n_{\max} = 3$ (the tail $\pi^2/6 - 49/36
\approx 0.284$ is not $\Delta$, nor any rational multiple of it);
and $\Delta$ is a boundary product at the truncation edge.
The combination rule $K = \pi(B + F - \Delta)$ therefore assembles
objects from three distinct tiers of the exchange-constant
taxonomy rather than specializing a single common Dirichlet
series, spectral construction, or arithmetic generator.

\paragraph{4. Negative results on candidate mechanisms.}
The three-tier reading is supported by the systematic elimination
of alternative unifying mechanisms.  Computation (exact rational
symbolic arithmetic and \texttt{mpmath} at 50 decimal places)
has ruled out, within the tolerance indicated:
(i) the Hopf-twist spectral comparison $S^3$ versus $S^1 \times S^2$
(cleanest near-miss $\sim 6.8 \times 10^{-3}$);
(ii) higher Casimir traces on $S^3$ at all orders through
$n_{\max} = 3$;
(iii) Hopf graph quotient Laplacian invariants for
$S^3 \to S^2$ over the $m$-fiber;
(iv) the discrete $S^1$ fiber spectral zeta (path-graph Laplacian,
in four weighting schemes);
(v) the continuous $S^1$ fiber via Mellin and heat-kernel expansions
(the Jacobi inversion $\theta_{S^1}(t) \sim \sqrt{\pi/t}$ forces every
order to be linear in $\pi$: the $\sqrt{\pi}\cdot\sqrt{\pi}\cdot\mathbb{Q}
= \pi\cdot\mathbb{Q}$ structure collapses any expected $\zeta_R(2)$
contribution to $\pi$ times a rational at every order);
(vi) $S^5$ spectral geometry, including zetas, Seeley--DeWitt
coefficients, and functional determinants (all rational, or
reducible to $\zeta_R(\mathrm{odd})$ and $\log 2$);
and (vii) a $\zeta$-combination scan for $\Delta$ over $3{,}228$
candidates built from Riemann and Hurwitz zetas, Bernoulli numbers,
and Stieltjes constants, with zero strict hits within $10^{-25}$
of $1/40$.

The structural reason for these negative results is that
round-sphere Laplace--Beltrami operators produce (i) integer powers
of $\pi$ in their volumes, (ii) rational Seeley--DeWitt
coefficients, and (iii) $\zeta_R(\mathrm{odd}) + \log 2$ combinations
in their spectral determinants---but never $\pi^2 = 2\zeta_R(2)$
additively next to a rational.  This is consistent with the
$\pi$-free certificate of the Bargmann--Segal lattice
(Paper~24), which shows that the
analogous lattice for the 3D harmonic oscillator on the Hardy
space of $S^5$ is bit-exactly rational in exact arithmetic; and the
same rigidity holds for $S^5$ Laplace--Beltrami spectral invariants
generally.

\paragraph{Reframed open question.}
The derivation question ``where does each piece of $K$ come from?''
has been partially answered: $B$ from a finite Casimir truncation
(with the $\kappa$--$B$ link of Eq.~\ref{eq:kappa_B_identity}), $F$
from an infinite Dirichlet series at the packing exponent
(Eq.~\ref{eq:F_dirichlet}), and $\Delta$ from a boundary product at
the truncation edge (irreducible to $\zeta$-content within the
candidate set tested).  The remaining structural question is no
longer a derivation question about any single component: it is a
question about why the \emph{additive} combination
$B + F - \Delta$, multiplied by the overall factor $\omega_2 = \pi$,
equals $\alpha^{-1}$ to $8.8 \times 10^{-8}$.  That is a question
about the conspiracy of three categorically different objects---a
finite sum, an infinite Dirichlet series, and a boundary
product---rather than the specialization of a single generator.
Paper~2's identification remains conjectural.


%% ========================================================================
%% VI. OBSERVABLE CLASSIFICATION
%% ========================================================================

\section{Observable Classification by Transcendental Content}
\label{sec:observable_classification}

The exchange constant taxonomy (Sec.~\ref{sec:taxonomy}) classifies the
\emph{constants themselves} by their mathematical origin.  A complementary
classification addresses the \emph{observables}: given a physical quantity
to be computed, what class of transcendental content does it require?

The conformal equivalence between the graph and continuous descriptions
(Paper~7) implies a natural partition of observables into three classes,
determined by what type of projection is needed to evaluate them.

\begin{enumerate}
  \item \textbf{Class~S (spectral observables):} Energy eigenvalues,
    degeneracies, selection rules, and transition amplitudes between
    eigenstates.  These require only the graph eigenvalue spectrum
    and the calibration constant $\kappa$.  No transcendental content
    beyond $\kappa = -1/16$ (rational) is needed.

  \item \textbf{Class~P (single-particle spatial observables):}
    Electron density $\rho(\mathbf{r})$, expectation values
    $\langle r^k \rangle$, momentum distributions, and any quantity
    requiring the wavefunction evaluated at a point in $\mathbb{R}^3$.
    These require the conformal projection $S^3 \to \mathbb{R}^3$,
    which introduces $\pi$ through the volume element
    $d\Omega_{S^3} = 2\pi^2$ and the spherical harmonic
    normalization $\int |Y_{lm}|^2 d\Omega = 1$ (with
    $d\Omega$ carrying $4\pi$).

  \item \textbf{Class~C (multi-particle correlated observables):}
    Correlation energies, adiabatic potential energy surfaces,
    multi-electron expectation values, and any quantity requiring
    the inter-particle coordinate.  These require the higher exchange
    constants: $e^a E_1(a)$ at Level~2 for $m \neq 0$ states
    (Sec.~\ref{sec:stieltjes}), $\mu(R)$ at Level~3
    (Sec.~\ref{sec:mu}), $\mu(\rho,R)$ at Level~4
    (Sec.~\ref{sec:mu_level4}), and the PK composition constant
    at Level~5.
\end{enumerate}

\noindent This classification is stated as a structural claim:

\medskip
\noindent\textbf{Claim~4.}
\textit{The class of transcendental content required to evaluate a
quantum-mechanical observable is determined by the type of projection
linking the graph description to the coordinate system in which the
observable is defined.  Spectral observables (Class~S) require no
transcendental content.  Single-particle spatial observables (Class~P)
require the Weyl constants ($\pi^{d/2}$).  Multi-particle correlated
observables (Class~C) require the embedding, flow, or composition
exchange constants of Sec.~\ref{sec:taxonomy}.}

\medskip

Claim~4 is falsifiable: for any observable, one can determine
its class, identify the required exchange constants, and verify
that no additional transcendental content enters the computation.
The following worked examples illustrate each class.

\subsection{Example 1: Hydrogen energy levels (Class~S)}

The hydrogen energy spectrum $E_n = -Z^2/(2n^2)$ is a spectral
observable.  The graph Laplacian on $S^3$ has eigenvalues
$\lambda_n = -(n^2 - 1)$, which are integers for all $n$
(Paper~7, Eq.~(12); Paper~0, Sec.~II).  The degeneracies
$g_n = n^2$ are integers.  The Fock conformal mapping (Paper~7,
Sec.~IV) converts these integer eigenvalues to the Rydberg
spectrum via the energy-shell constraint $p_0^2 = -2E$
and the stereographic focal length:
\begin{equation}
  E_n = -\frac{Z^2}{2n^2}\;\text{Ha},
  \label{eq:hydrogen_class_s}
\end{equation}
with the calibration constant $\kappa = -1/16$ (rational)
mediating between graph eigenvalues and physical energies
(Sec.~\ref{sec:kappa_derivation}).  Every quantity in this
derivation---eigenvalues, degeneracies, $\kappa$---is an
integer or rational number.  No transcendental content enters.

Selection rules provide a second example.  The dipole selection
rule $\Delta l = \pm 1$ follows from the triangle inequality
on the Gaunt integral
$\int Y_{l_1 m_1} Y_{1q} Y_{l_2 m_2}\,d\Omega$, which
reduces to a product of Wigner $3j$ symbols---rational numbers
computed from factorials of integers~\cite{loutey_paper0}.
No transcendental content enters.

\subsection{Example 2: Electron density (Class~P)}

The electron density $\rho(\mathbf{r}) = |\psi(\mathbf{r})|^2$
requires evaluating the wavefunction in position space.  On the
graph, the state $|n, l, m\rangle$ is a node with an integer
label; in $\mathbb{R}^3$, the corresponding wavefunction is
$\psi_{nlm}(\mathbf{r}) = R_{nl}(r)\,Y_{lm}(\theta,\phi)$.
The conformal projection $S^3 \to \mathbb{R}^3$ (Paper~7,
Sec.~III) introduces the conformal factor
$\Omega = 2p_0/(p^2 + p_0^2)$, which mediates between the
unit-sphere measure $d\Omega_{S^3}$ and the flat-space volume
element $d^3\mathbf{r}$.  The spherical harmonic normalization
$\int |Y_{lm}|^2\,d\Omega = 1$ introduces $\pi$ through the
solid angle $\oint d\Omega = 4\pi$.  The radial normalization
$\int_0^\infty |R_{nl}|^2\,r^2\,dr = 1$ introduces no additional
transcendentals for hydrogen (the integrals are rational multiples
of $a_0^{-3}$ times exponentials, evaluated in closed form).

The density is thus a Class~P observable: it requires the $\pi$
that enters through the projection from the dimensionless $S^3$
topology to the dimensionful $\mathbb{R}^3$ coordinate system.

\subsection{Example 3: Helium ground state energy (Class~C)}

The helium ground state energy $E_0 = -2.9037$~Ha~\cite{drake2006}
is a correlated observable.  The Level~3 hyperspherical solver
(Paper~13) computes it via the adiabatic potential $U(R)$, which
depends on the angular eigenvalue $\mu(R)$.  The function $\mu(R)$
satisfies the characteristic polynomial $P(R,\mu) = 0$
(Sec.~\ref{sec:algebraic_curve}), with coefficients in
$\mathbb{Q}(\pi, \sqrt{2})$ inherited from the angular coupling
matrix~\cite{loutey_track_p1}.  The energy is obtained by
integrating the hyperradial Schr\"odinger equation with
$U(R)$ as the effective potential---a definite integral over
the algebraic curve, producing a transcendental number.

On the $S^3 \times S^3$ product space (Level~3 FCI), the same
energy can in principle be computed without $\mu(R)$: the
Hamiltonian matrix has entries given by Gaunt integrals (rational)
and the Slater $F^0$ direct Coulomb integrals (algebraic on $S^3$,
Paper~7 Sec.~VI).  The FCI result (0.35\% error) demonstrates that
the correlated energy \emph{can} be approximated within the
algebraic framework, but the $\mu(R)$ parameterization is needed
to achieve sub-0.1\% accuracy.  The transition from Class~S to
Class~C reflects the physical content: electron-electron
correlation couples the two $S^3$ factors, and the coupling
strength varies with the hyperradial coordinate $R$---a spatial
degree of freedom that exists only in the projected description.

\paragraph{Three-layer decomposition (v2.6.0).}
The algebraic Casimir CI (Track~DI Sprint~2) makes the exchange
constant structure of the He ground state fully explicit.  At each
$n_{\max}$, the FCI matrix $H(k) = Bk + Ck^2$ has rational entries
as a function of the orbital exponent~$k$, so the variationally
optimal energy $E^*$ is an algebraic number---a root of a polynomial
with rational coefficients.  At $n_{\max} = 1$: $k^* = 27/16$,
$E^* = -729/256$~Ha (exact rationals).  This reveals a three-layer
structure:
\begin{center}
\begin{tabular}{lll}
\hline
Layer & Content & Exchange type \\
\hline
Rational & Slater integrals, Gaunt coefficients & None (graph-intrinsic) \\
Algebraic & Optimal exponents at finite $n_{\max}$ & Calibration \\
Transcendental & e-e cusp ($\alpha = \pi/4$ on $S^5$) & Embedding \\
\hline
\end{tabular}
\end{center}
The Fock energy-shell self-consistency $k^2 = -2E$, which is exact
at Level~1 (one electron, one scale), over-constrains the
two-electron problem: self-consistent~$k$ gives 5--13\% error vs.\
1.5--1.9\% for variational~$k$.  This is the first instance where a
single calibration constant is insufficient---the transition from
calibration to embedding exchange constants corresponds to the
physical content of electron-electron correlation.

\subsection{Boundary cases}

Two boundary cases clarify the classification:

\begin{itemize}
  \item \textbf{Ionization energies} are spectral (Class~S) when
    computed as eigenvalue differences $\Delta E = E_n - E_m$,
    which involve only $\kappa$ and integer quantum numbers.
    However, photoionization \emph{cross sections} require the
    dipole matrix element $\langle \psi_f | \mathbf{r} | \psi_i
    \rangle$, which is a spatial observable (Class~P) because
    it evaluates the wavefunction at positions in $\mathbb{R}^3$.

  \item \textbf{Two-electron energies} computed via FCI on
    $S^3 \times S^3$ are formally Class~S (graph eigenvalues
    of the product-space Hamiltonian), but achieve limited
    accuracy due to the product-basis completeness deficiency
    (Sec.~\ref{sec:mu_level4}).  High-accuracy two-electron
    energies require the Class~C exchange constants.
    The classification thus tracks both the type of observable
    and the accuracy target.
\end{itemize}

The classification does not assert that one class is more
fundamental than another.  The conformal equivalence
(Paper~7) ensures that all three descriptions---graph, projected
single-particle, correlated multi-particle---are mathematically
equivalent.  The classification identifies which exchange
constants mediate between them.


%% ========================================================================
%% VII. DISCUSSION
%% ========================================================================

\section{Discussion}
\label{sec:discussion}

\subsection{Transcendence and the graph design principle}

Exchange constants are transcendental because they convert between
incommensurable descriptions: a discrete count (eigenvalues,
representation labels) and a continuous measure (volume, curvature,
geodesic length).  Weyl constants carry $\pi^{d/2}$, and $\pi$ is
transcendental by Lindemann's theorem.  Not all functions in the
hierarchy are transcendental: $\mu(R)$ satisfies a polynomial $P(R,\mu) = 0$
(Sec.~\ref{sec:algebraic_curve}), making it algebraic in~$R$, with
transcendental content inherited from angular coupling
coefficients~\cite{berger2003}.  The irreducible transcendentals are
the \emph{constants} that enter the coupling matrix, not the
parameterized functions built from them.

When the computation stays on the graph, exchange constants are
unnecessary.  This suggests a design principle: \textit{minimize
coordinate projections} to minimize transcendental content.  The
ideal solver stays on the graph; every departure incurs a
transcendental toll.

\subsection{Exchange constants are computationally cheap}

Conceptual complexity does not correlate with computational cost:
$\kappa$ is one multiplication; $e^a E_1(a)$ is one function
evaluation seeding an $O(N^2)$ recurrence; $\mu(R)$ is a
$10$--$30$ dimensional matrix diagonalization per $R$-point.
The hierarchy tracks curvature mismatch severity, not computational
difficulty.  At Level~2, identifying the Stieltjes seed and absorbing
it into the basis \emph{reduces} cost (eliminating quadrature) while
improving convergence (Sec.~\ref{sec:stieltjes}).

\subsection{The hierarchy of transcendence}

Exchange constants increase in complexity with curvature mismatch:
constant ($\kappa = -1/16$, Level~1) $\to$ single transcendental
seed ($e^a E_1(a)$, Level~2) $\to$ algebraic function of one
parameter ($\mu(R)$ satisfying $P(R,\mu) = 0$ over
$\mathbb{Q}(\pi,\sqrt{2})$, Level~3) $\to$ piecewise-algebraic
function of two parameters ($\mu(\rho,R)$, Level~4).  This mirrors
the mathematical hierarchy: Weyl constants are periods
($\pi = \int dx/(1+x^2)$), Stieltjes seeds are exponential periods
($E_1(a) = \int_1^\infty e^{-at}/t\,dt$), and adiabatic eigenvalues
are algebraic functions of the continuous parameter built from
transcendental coefficients.


\subsection{Structural parallels}

The pattern parallels string compactification, where compact internal
geometry determines low-energy couplings~\cite{polchinski1998}.
We note also that
$\mathrm{SO}(4,2) = \mathrm{Iso}(\mathrm{AdS}_5)$, the isometry group
of five-dimensional anti-de~Sitter space, is the conformal group of the
boundary $\mathbb{R}^{3,1}$; the Fock projection
$S^3 \to \mathbb{R}^3$ is a Euclidean section of this conformal
boundary correspondence.

\subsection{The algebraic curve refinement}
\label{sec:algebraic_curve}

The characteristic polynomial analysis~\cite{loutey_track_p1}
refines the exchange constant taxonomy by distinguishing between
transcendental \emph{constants} and algebraic \emph{functions} of
continuous parameters:
\begin{itemize}
  \item \textbf{Truly irreducible transcendentals:}
    The embedding constants $e^a E_1(a)$ (Level~2, $m \neq 0$) and the
    angular coupling integrals $c/\pi$ and $c\sqrt{2}/\pi$
    (Level~3 and~4) are irreducible---they cannot be eliminated by
    any algebraic manipulation.  These are the genuine ``exchange
    constants'' in the Weyl--Selberg sense.
  \item \textbf{Algebraic functions of transcendental data:}
    The adiabatic eigenvalues $\mu(R)$ and $\mu(\rho,R)$ are
    algebraic functions whose coefficients inherit transcendentals
    from the coupling matrix.  The $R$-dependence (or $(\rho,R)$-dependence)
    is algebraic at every truncation: $\mu$ satisfies a polynomial
    equation $P(R,\mu) = 0$ of degree $d$ equal to the matrix dimension.
    At $l_{\max} = 0$, $P$ is linear and $\mu(R)$ is an explicit
    rational function of $R$ (with transcendental coefficients).
    At $l_{\max} = L$, $P$ has degree $L+1$ in both $R$ and $\mu$.
  \item \textbf{Calibration constant:} $\kappa = -1/16$ is rational
    (Sec.~\ref{sec:kappa_derivation}).
\end{itemize}

This refinement has a practical implication: the ``transcendental toll''
of the adiabatic parameterization is \emph{inherited}, not \emph{accrued}.
Moving from the product space $S^3 \times S^3$ to hyperspherical
coordinates does not introduce new transcendental content; it
rearranges existing transcendentals (from the angular coupling)
into an algebraic function of $R$.  The toll is paid once, in the
coupling matrix, and the $R$-dependence is algebraically determined.

This does not yet eliminate the computational cost of pointwise
diagonalization: at $l_{\max} \geq 1$, the characteristic polynomial
has degree $\geq 2$ in $\mu$ and must be solved numerically at each
$R$-point.  But it does mean that the solver operates on a known
polynomial, not an opaque transcendental function.  Whether this
algebraic structure can be exploited computationally---for instance,
via resultant methods, homotopy continuation, or algebraic
PES representations---is an open question.

This algebraic structure is specific to Level~3, where the angular
Hamiltonian $H(R) = H_0 + R \cdot V_C$ is a linear matrix pencil.
At Level~4, the split-region Legendre expansion of the nuclear charge
function introduces piecewise-smooth $\rho$-dependence that breaks
the pencil structure (Track~S).  The Level~4 angular eigenvalues
are piecewise-algebraic but do not satisfy a global characteristic
polynomial.

\subsection{Derivation of $\kappa = -1/16$}
\label{sec:kappa_derivation}

The original identification of $\kappa = -1/16$ (Paper~0) was
empirical: the unique constant that maps GeoVac graph eigenvalues
onto the Rydberg energy levels.  We now derive this value and
classify its origin.

\textit{Graph Laplacian spectrum.}  The GeoVac lattice
(Paper~0) discretizes $S^3$ with states $|n,l,m\rangle$
connected by angular transitions $m \leftrightarrow m \pm 1$
and radial transitions $n \leftrightarrow n \pm 1$.  Interior
nodes have degree $d_{\max} = 4$ (two angular, two radial
neighbors).  The graph Laplacian $\mathcal{L} = D - A$ is
positive semidefinite; its maximum eigenvalue satisfies
$\lambda_{\max} \leq 2 d_{\max} = 8$, with the bound saturated
in the continuum limit ($n_{\max} \to \infty$).  Numerical
verification confirms monotonic convergence:
$\lambda_{\max}(n_{\max}{=}5) = 6.62$,
$\lambda_{\max}(10) = 7.61$,
$\lambda_{\max}(20) = 7.90 \to 8$.

\textit{Energy calibration.}  The graph Hamiltonian
$H = \kappa \mathcal{L}$ (no node weights) has its ground
state at energy $\kappa \lambda_{\max}$, since $\kappa < 0$
maps the largest graph eigenvalue to the lowest physical energy.
Matching this to the hydrogen ground state:
\begin{equation}
  \kappa \lambda_{\max} = E_H = -\tfrac{1}{2}
  \implies
  \kappa = \frac{-1/2}{8} = -\frac{1}{16}.
  \label{eq:kappa_derived}
\end{equation}

\textit{Factor identification.}  The value $-1/16 = -(1/2)(1/8)$
decomposes into:
\begin{enumerate}
  \item $1/2$: the non-relativistic kinetic energy normalization
    $E = p^2/(2m)$, or equivalently the energy-shell constraint
    $E = -p_0^2/2$;
  \item $1/8 = 1/(2 d_{\max})$: the inverse of the maximum
    graph Laplacian eigenvalue, determined by the vertex degree
    of the lattice discretization.
\end{enumerate}

\textit{Heat kernel test.}  The Seeley--DeWitt coefficients on unit
$S^3$ ($a_0 = 2\pi^2$, $a_1 = 2\pi^2$, $a_2 = \pi^2$) produce no
ratio matching $1/16$; neither does the Weyl counting function
$N(\Lambda) \sim \Lambda^{3/2}/3$ or the spectral zeta.  The conformal
coupling $c_3 = 1/8$ for $d=3$ numerically equals $1/\lambda_{\max}$,
but this is coincidental: $c_3$ is a Yamabe invariant while
$1/\lambda_{\max}$ depends on graph vertex degree.

\textit{Fock weight derivation (Note added v2.26.1).}  The same value
is derivable independently from the Fock projection itself, without
reference to graph vertex degree.  The $l = 0$ eigenfunctions on $S^3$
are Gegenbauer polynomials $C^1_{n-1}(\cos\chi)$ (equivalently,
Chebyshev polynomials of the second kind $U_{n-1}(\cos\chi)$).  Their
three-term recurrence gives transition amplitudes of $1/4$ between
adjacent shells, so the squared Fock coupling is
\begin{equation}
  c^2(n,l) = \tfrac{1}{16}\bigl[1 - l(l{+}1)/(n(n{+}1))\bigr],
  \label{eq:fock_coupling_p18}
\end{equation}
with $c^2(n,0) = 1/16$ universally for all $n$.  Equivalently,
$1/16 = 1/\Omega^4(0)$ where $\Omega(0) = 2$ is the stereographic
conformal factor at $p = 0$.  The graph-degree derivation
(Eq.~\ref{eq:kappa_derived}) and the Fock weight derivation converge
on the same value because the graph's maximum vertex degree
$d_{\max} = 4$ and the conformal factor $\Omega(0) = 2$ are
independently determined by the same underlying $S^3$ geometry:
$2 d_{\max} = \Omega(0)^4 = 2^4/2 = 8$.

A notable special case: $c^2(4,3) = 1/40 = \Delta$ (Paper~2), linking
the highest-$l$ Fock coupling at the Paper~2 cutoff to the Dirac
degeneracy $g_3^{\mathrm{Dirac}} = 40$.

\textit{Classification.}  $\kappa$ is reclassified from ``calibration''
to ``conformal'' in Table~\ref{tab:taxonomy}: its numerical value
$1/16$ is derivable from the continuous conformal geometry of the Fock
projection ($1/\Omega^4(0)$), not merely from the discrete graph
connectivity.  The two derivations (graph degree and conformal factor)
agree because the graph discretizes $S^3$, but the conformal-factor
reading is more fundamental---it would hold for any discretization of
$S^3$, not only the GeoVac lattice.  The Z-universality noted in
Paper~0 ($\kappa \to \kappa Z^2$ for any nuclear charge)
reflects the Z-independence of the graph structure, with only
the energy-shell scaling $E \propto Z^2$.


%% ========================================================================
%% VIII. CONCLUSION
%% ========================================================================

\subsection{Structural incommensurability within GeoVac observables:
a meta-pattern}
\label{sec:meta_observation}

The exchange-constant taxonomy developed in the preceding sections
has repeatedly encountered a distinctive pattern: single framework
observables in GeoVac, when unpacked, resolve into sums or
products of categorically distinct transcendental or spectral
objects with \emph{no common generator}. The pattern has now been
observed in three independently-derived three-axis classifications,
and we note it here as a meta-observation that may eventually
rise to theorem status.

\paragraph{Instance~1 (Paper~28 QED on $S^3$):}  Perturbative QED
sums on $S^3$ produce three transcendental classes simultaneously:
even-zeta content ($\pi^{\mathrm{even}}$, from second-order
Riemannian operators), odd-zeta content ($\zeta(3), \zeta(5)$,
from first-order Dirac), and Dirichlet-$L$ content ($G = \beta(2),
\beta(4)$, from vertex topology). These three classes appear
together in the two-loop sunset sum but are discriminated by a
three-axis classification (operator order $\times$ bundle type
$\times$ vertex parity), and no single mechanism generates all
three. \cite{loutey_paper28}

\paragraph{Instance~2 (Phase-4G $\alpha$ decomposition):}  The
conjectural combination rule $K = \pi(B + F - \Delta)$ has three
ingredients with categorically distinct origins: $B = 42$ is a
finite Casimir trace on the scalar Laplacian at $m = 3$, $F = \pi^2/6$
is the infinite Dirichlet series of Fock degeneracies at the
packing exponent $d_{\max} = 4$, and $\Delta^{-1} = 40$ is the
degeneracy of the third Dirac mode $g_3^{\mathrm{Dirac}} =
2 \cdot 4 \cdot 5$ on the \emph{spinor} bundle. Sprint 4 analysis
confirmed that no common generator produces all three: $K$ is an
additive combination across three structurally-independent
spectral objects. (Phase~4G memo,
\texttt{debug/data/track\_alpha\_phase4g/track\_k\_analysis.md}.)

\paragraph{Instance~3 (RH-sprint Weil dictionary):}  The Ihara
zeta side (combinatorial, closed non-backtracking walks) and the
spectral Dirichlet side (Dirac-eigenvalue Dirichlet series) give
\emph{opposite} verdicts on GeoVac. Zero spacing: Ihara Poisson
($\mathrm{CV} \approx 1$), spectral GUE ($\mathrm{CV} \approx
0.35$). Dirichlet-$L$ content: Ihara absent (bipartiteness makes
the $\chi_{-4}$ support disjoint from walk-length support); spectral
present (the closed-form Eq.~(\ref{eq:rh_j_identity})). Hilbert--P\'olya
compatibility: Ihara categorically excluded; spectral compatible
at current sample precision. The two sides are not two views of
the same object; they are two distinct framework observables
computed from the same Dirac-$S^3$ data, with structural support
that is disjoint in the bipartite-$\chi_{-4}$ sense.

\paragraph{The meta-pattern.}  Across these three instances, the
framework keeps finding that a single-observable statement
("$\zeta$-zero structure on GeoVac", "$K = \pi(B + F - \Delta)$",
"two-loop QED spectral sum on $S^3$") resolves under structural
analysis into a sum or superposition of objects that do not share
a common structural generator. The appropriate taxonomy is the
three-axis-classification style developed here; the objects at
each cell are what they are, not derivations of each other.

This is consistent with --- and arguably a consequence of --- the
Coulomb/HO asymmetry established in Paper~24~\cite{loutey_paper24}:
the graph is rational/integer on both bundles, but the projection
structures (second-order Riemannian for the scalar, first-order
complex-analytic for the HO holomorphic sector, Dirac for the
spinor) are categorically different, and their transcendental
projections onto physical coordinates do not interpolate.
We name the pattern for future reference --- \emph{structural
incommensurability within GeoVac observables} --- and note that
this is an \emph{empirical} observation at the level of the three
documented instances, not a theorem. Its status may sharpen if
further framework observables decompose the same way or if a
bridge is found that fails cleanly.

\paragraph{Operational consequence.}  For taxonomy construction
(Sec.~\ref{sec:taxonomy}), the consequence is that multi-component
observables should be decomposed and each component classified
separately, rather than the whole observable forced into a single
tier. This changes nothing about Sec.~\ref{sec:observable_classification}'s
S/P/C classification of atomic observables (which are fine as
they stand at the observable level) but it does mean that
composite structural results, such as Paper~2's $K$ or Paper~28's
two-loop sunset, should be broken into their components before
taxonomic assignment.

% ------------------------------------------------------------------

\section{Conclusion}
\label{sec:conclusion}

This paper establishes three results:

\paragraph{1. The taxonomy.}
Transcendental conversion factors in the GeoVac hierarchy are
classified by the type of compactness mismatch they mediate:
intrinsic ($\pi^{d/2}$), calibration ($\kappa = -1/16$),
embedding ($e^a E_1(a)$, $1/r_{12}$), flow ($\mu(R)$,
$\mu(\rho,R)$), and composition (PK pseudopotential).  The
Level~3 adiabatic eigenvalue $\mu(R)$ is an algebraic function
of~$R$ with transcendental coefficients; the Level~4 eigenvalue
$\mu(\rho,R)$ is piecewise-algebraic.  All exchange constants
vanish when the computation stays on the graph.

\paragraph{2. The three-axis classification and flat-space consistency.}
Operator order on the compact space determines the motivic weight
of the transcendental content: second-order operators
($\Delta_{\text{LB}}$, $D^2$) produce $\pi^{\text{even}}$
exclusively (Eq.~\ref{eq:zeta_d2}); first-order operators ($D$)
produce odd-zeta content ($\zeta_R(3)$, Eq.~\ref{eq:dirac_dirichlet_zeta3}).
Bundle type modulates coefficients; diagram topology introduces
Dirichlet $L$-function values ($G$, $\beta(4)$) via the vertex
parity selection rule (Eqs.~\ref{eq:d_even4}--\ref{eq:d_odd4}).
The three-axis taxonomy is validated as $R$-independent: the ratio
$D(s,R)/R^s$ is a structural invariant, and Weyl's law confirms the
spectral-to-momentum correspondence in the flat-space limit.
The spinor-intrinsic ring $R_{\mathrm{sp}}$ (Eq.~\ref{eq:ring_sp})
captures all relativistic content from Breit--Pauli through
full-Dirac corrections.

\paragraph{3. The QED beta function from spectral data.}
The operator-order discriminant predicts: no odd-zeta at one loop
(verified: the Schwinger~\cite{schwinger1948} result
$a_e^{(1)} = \alpha/(2\pi)$ involves only $\pi$); $\zeta_R(3)$ at
two loops (verified: Petermann~\cite{petermann1957}).  The Seeley--DeWitt
coefficients $a_0, a_1, a_2$ of $D^2$ on $S^3$, which control the
one-loop vacuum polarization, are purely $\pi^{\text{even}}$ by
Eq.~(\ref{eq:zeta_d2}).

\paragraph{Verified and open.}
The taxonomy has been validated across the full GeoVac hierarchy
through Level~5.  Claim~4 (transcendental content determined by
projection type) has held across all tested cases.
What remains open: (i)~whether the three-axis taxonomy predicts
the physical two-loop coefficient (not just the transcendental class);
(ii)~the conjectural combination rule $K = \pi(B + F - \Delta)$ of
Paper~2, whose three-tier structure (Sec.~\ref{sec:alpha_three_tier})
is fully characterized but whose additive conspiracy remains
unexplained.

\paragraph{Design principle.}
Staying on the graph eliminates exchange constants by eliminating
projections.  The program's effort to algebraicize solvers (Paper~12's
Neumann expansion, Laguerre recurrences, Gaunt replacements,
Paper~19's cross-center $V_{ne}$ via incomplete gamma functions) is,
in retrospect, an effort to eliminate exchange constants.  The algebraic
curve result (Sec.~\ref{sec:algebraic_curve}) shows this has been more
successful than previously recognized.

\paragraph{Note added v2.6.0.}
The He exchange constant census (Track~DI) validates the taxonomy
computationally.  Graph-native FCI achieves 0.19\% at $n_{\max} = 7$
with zero free parameters.  Three-layer structure: (1)~rational
(graph $h_1$ + Slater $V_{ee}$, 98.6\%); (2)~topological (graph
off-diagonal $T_\pm$, calibration $\kappa$, dominant); (3)~transcendental
(e-e cusp, embedding, ${\sim}0.14\%$ floor).  FCI basis invariance
confirmed: the floor is embedding content (cusp), not basis mismatch.

\paragraph{Note added v2.26.1 (Seeley--DeWitt cancellation).}
On unit $S^{3}$, the Seeley--DeWitt coefficient ratio
$a_{k}^{D^{2}} / a_{k}^{\Delta_{\mathrm{LB}}} = 4 = \dim(S)$
at every order $k$, where $S$ denotes the spinor bundle.  This ratio
is exact: the asymptotic mode density of the Camporesi--Higuchi Dirac
spectrum is $g_{n}^{\mathrm{Dirac}} \sim 4n^{2}$, four times the
scalar density $n^{2}$, and the SD expansion depends only on this
asymptotic density.  As a consequence, the perturbative
Connes--Chamseddine supertrace
$\mathrm{Str}[f(D^{2}/\Lambda^{2})] \equiv
 \mathrm{Tr}_{S}[f] - \tfrac{1}{4}\mathrm{Tr}_{D}[f]$
vanishes identically in the asymptotic expansion.  The three
ingredients of $K/\pi = B + F - \Delta$ are all invisible to the
SD expansion: $B$ is a finite Casimir trace, $F$ is an arithmetic
Dirichlet series, and $\Delta^{-1} = 40$ is the Euler--Maclaurin
upper boundary term of the Dirac mode-count sum at $n_{\mathrm{CH}}
= 3$.  The sign on $\Delta$ in the combination rule $K = \pi(B + F -
\Delta)$ is the standard $(-1)^{F}$ grading: $B$ and $F$ are
scalar-sector (bosonic) quantities that enter with $(+)$, while
$\Delta$ is a spinor-sector (fermionic) quantity that enters with
$(-)$.  $K/\pi$ lives entirely in the non-perturbative sector of the
spectral-action supertrace on~$S^{3}$.

\begin{acknowledgments}
Computational implementation and documentation drafting were
performed with the assistance of large language models (Anthropic
Claude) under the author's scientific direction.  All results
were validated against analytical benchmarks.
\end{acknowledgments}


%% ========================================================================
%% BIBLIOGRAPHY
%% ========================================================================

\begin{thebibliography}{30}

\bibitem{loutey_paper0}
J.~Loutey,
``The Geometric Atom: Quantum State Space as a Packing Problem,''
GeoVac Paper~0 (2026).

\bibitem{loutey_paper1}
J.~Loutey,
``The Geometric Atom: Quantum Mechanics as a Packing Problem,''
GeoVac Paper~1 (2026).

\bibitem{loutey_paper7}
J.~Loutey,
``The Dimensionless Vacuum: Recovering the Schr\"odinger Equation
from Scale-Invariant Graph Topology,''
GeoVac Paper~7 (2026).

\bibitem{loutey_paper11}
J.~Loutey,
``The Molecular Fock Projection: Diatomic Molecules as Prolate
Spheroidal Lattices,''
GeoVac Paper~11 (2026).

\bibitem{loutey_paper13}
J.~Loutey,
``The Hyperspherical Lattice: Two-Electron Atoms as Coupled
Channel Graphs,''
GeoVac Paper~13 (2026).

\bibitem{loutey_fci}
J.~Loutey,
``Full configuration interaction on a sparse graph Hamiltonian:
Multi-electron atoms via relational lattice indexing,''
GeoVac FCI Paper (2026).

\bibitem{loutey_paper2}
J.~Loutey,
``The Fine Structure Constant from Spectral Geometry of the Hopf
Fibration,''
GeoVac Paper~2 (2026).

\bibitem{loutey_paper24}
J.~Loutey,
``The Bargmann--Segal Lattice: A $\pi$-Free Discretization of the
3D Harmonic Oscillator on $S^5$,''
GeoVac Paper~24 (2026).

\bibitem{loutey_paper28}
J.~Loutey,
``QED on $S^3$: Spectral Geometry of Perturbative Transcendentals,''
GeoVac Paper~28 (2026).

\bibitem{loutey_paper29}
J.~Loutey,
``The GeoVac Hopf Graphs Are Ramanujan: Ihara Zeta, Graph-RH, and
a Scope Boundary on Selberg-on-Hydrogen,''
GeoVac Paper~29 (2026).

\bibitem{loutey_track_j}
J.~Loutey,
``Track J: Algebraic Level 2 Radial Solver for $m \neq 0$ States,''
GeoVac supplementary material (2026).

\bibitem{loutey_track_g}
J.~Loutey,
``Track G: Perturbation Series for $\mu(R)$,''
GeoVac supplementary material (2026).

\bibitem{loutey_track_m}
J.~Loutey,
``Track M: Prolate CI $l_{\max}$ Convergence Study,''
GeoVac supplementary material (2026).

\bibitem{loutey_track_p1}
J.~Loutey,
``Track P1: Algebraic Curve for $\mu(R)$---Characteristic Polynomial
of the Level~3 Angular Hamiltonian,''
GeoVac supplementary material (2026).

\bibitem{loutey_track_s}
J.~Loutey,
``Track S: Level~4 Algebraic Curve---Structural Negative Result,''
GeoVac supplementary material (2026).

\bibitem{drake2006}
G.~W.~F.~Drake,
``High precision calculations for helium,''
in \textit{Springer Handbook of Atomic, Molecular, and Optical Physics},
edited by G.~W.~F.~Drake (Springer, New York, 2006), pp.~199--219.

\bibitem{dlmf_laguerre}
NIST Digital Library of Mathematical Functions,
\S18.9.13, \url{https://dlmf.nist.gov/18.9.13}.

\bibitem{weyl1911}
H.~Weyl,
``\"Uber die asymptotische Verteilung der Eigenwerte,''
Nachr.\ K\"onigl.\ Ges.\ Wiss.\ G\"ottingen, Math.-phys.\ Kl.\
(1911), pp.\ 110--117.

\bibitem{selberg1956}
A.~Selberg,
``Harmonic analysis and discontinuous groups in weakly symmetric
Riemannian spaces with applications to Dirichlet series,''
J.\ Indian Math.\ Soc.\ \textbf{20}, 47--87 (1956).

\bibitem{berger2003}
M.~Berger,
\textit{A Panoramic View of Riemannian Geometry}
(Springer, Berlin, 2003).

\bibitem{lichnerowicz1963}
A.~Lichnerowicz,
\textit{Spineurs harmoniques},
C.~R.\ Acad.\ Sci.\ Paris \textbf{257}, 7--9 (1963).

\bibitem{camporesi_higuchi1996}
R.~Camporesi and A.~Higuchi,
On the eigenfunctions of the Dirac operator on spheres and
real hyperbolic spaces,
\textit{J.\ Geom.\ Phys.}\ \textbf{20}, 1--18 (1996).

\bibitem{peter_weyl1927}
F.~Peter and H.~Weyl,
``Die Vollst\"andigkeit der primitiven Darstellungen einer
geschlossenen kontinuierlichen Gruppe,''
\textit{Math.\ Ann.}\ \textbf{97}, 737--755 (1927).

\bibitem{polchinski1998}
J.~Polchinski,
\textit{String Theory}, Vol.~2: Superstring Theory and Beyond
(Cambridge University Press, Cambridge, 1998).

\bibitem{goncharov1999}
A.~B.~Goncharov,
``Volumes of hyperbolic manifolds and mixed Tate motives,''
\textit{J. Amer. Math. Soc.} \textbf{12}, 569--618 (1999).

\bibitem{zagier1994}
D.~Zagier,
``Values of zeta functions and their applications,''
in \textit{First European Congress of Mathematics}, Vol.~II
(Birkh\"auser, Basel, 1994), pp.~497--512.

\bibitem{schwinger1948}
J.~Schwinger,
``On quantum-electrodynamics and the magnetic moment of the electron,''
\textit{Phys.\ Rev.}\ \textbf{73}, 416--417 (1948).

\bibitem{petermann1957}
A.~Petermann,
``Fourth order magnetic moment of the electron,''
\textit{Helv.\ Phys.\ Acta} \textbf{30}, 407--408 (1957).

\bibitem{sommerfield1957}
C.~M.~Sommerfield,
``Magnetic dipole moment of the electron,''
\textit{Phys.\ Rev.}\ \textbf{107}, 328--329 (1957).

\bibitem{adamchik2003}
V.~S.~Adamchik,
``A certain series associated with Catalan's constant,''
\textit{Z.\ Anal.\ Anwend.}\ \textbf{22}, 817--826 (2003).

\bibitem{seeley1967}
R.~T.~Seeley,
``Complex powers of an elliptic operator,''
\textit{Proc.\ Sympos.\ Pure Math.}\ \textbf{10}, 288--307 (1967).

\bibitem{minakshisundaram1949}
S.~Minakshisundaram and {\AA}.~Pleijel,
``Some properties of the eigenfunctions of the Laplace-operator
on Riemannian manifolds,''
\textit{Canad.\ J.\ Math.}\ \textbf{1}, 242--256 (1949).

\bibitem{gilkey1975}
P.~B.~Gilkey,
``The spectral geometry of a Riemannian manifold,''
\textit{J.\ Differential Geom.}\ \textbf{10}, 601--618 (1975).

\end{thebibliography}

\end{document}
